% Generated mechanically from frozen input inputs/p09/ja/src/spaces-groupoids.tex.
% Do not edit this generated file; edit the bound locale source instead.

% OK, start here.
%

\setcounter{chapter}{77}
\chapter{代数空間の亜群}



\phantomsection
\label{spaces-groupoids-section-phantom}


\section{はじめに}
\label{spaces-groupoids-section-introduction}

\noindent
本章では、代数空間の亜群に関する一般論を扱う。Keel と Mori による美しい論文
\cite{K-M} を読むことを勧める。

\medskip\noindent
ここで述べることの多くは、亜群スキームの章で述べたことの繰り返しである。
Groupoids, Section \ref{groupoids-section-introduction} を参照せよ。商スタックに
関する議論は本章で新たに加わる。


\section{規約}
\label{spaces-groupoids-section-conventions}

\noindent
すべてのスキームは大 fppf サイト $\Sch_{fppf}$ に含まれるものと常に仮定する。
また、考察するすべての環 $A$ は、$\Spec(A)$ がこの大サイトの対象と同型である
という性質をもつものとする。

\medskip\noindent
$S$ をスキームとし、$X$ を $S$ 上の代数空間とする。本章および次章では、
$X\times_S X$ と書く。これは $X$ とそれ自身との積（$S$ 上の代数空間の圏における
積）であり、$X\times X$ とは書かない。

\medskip\noindent
射影の添字を $0$ から始めるという規約を引き続き用いる。したがって
$\text{pr}_0:X\times_S Y\to X$ および
$\text{pr}_1:X\times_S Y\to Y$ である。




\section{記法}
\label{spaces-groupoids-section-notation}

\noindent
$S$ をスキームとする。これを基礎スキームとし、すべての代数空間は $S$ 上にあるものとする。
$B$ を $S$ 上の代数空間とする。これを基礎代数空間とし、しばしば他の代数空間および
スキームも $B$ 上にあるものとする。$X$ が $B$ 上の代数空間であるというとき、それは
$X$ が $S$ 上の代数空間で、構造射 $X\to B$ を備えるという意味である。また、文字 $T$ は
$B$ 上の「試験」スキームを表すために使うよう努める。言い換えると、$T$ は構造射
$T\to B$ を備えたスキームである。この状況では $X(T)$ を、$T$-値点の集合であって
$X$ のものを $B$ 上で考えたものとして表す。式で書けば、
$$
X(T)=\Mor_B(T,X).
$$
同様に、第2の代数空間 $Y$ が $B$ 上に与えられたとき、
$$
X(Y)=\Mor_B(Y,X).
$$
とおく。代数空間 $X$、$Y$ が上のように $B$ 上にあり、射 $f:X\to Y$ が
$B$ 上に与えられているとする。任意のスキーム $T$ で $B$ 上にあるものに対し、集合の誘導写像
$$
f:X(T)\longrightarrow Y(T)
$$
を得る。これはスキーム $T$ に関して関手的であり、そのスキームは $B$ 上にある。$f$ は
$(\Sch/S)_{fppf}$ 上の層の写像で、層 $B$ 上にあるので、この規則を定め、
また $f$ はこの規則によって定まることは明らかである。より一般に、ファイバー積間の
写像にも同じ記法を用いる。例えば、$X$、$Y$、$Z$ が $B$ 上の代数空間であり、
$m:X\times_B Y\to Z\times_B Z$ が $B$ 上の代数空間の射ならば、$m$ は
$T$-値点間の写像の集まり
$$
X(T)\times Y(T)\longrightarrow Z(T)\times Z(T).
$$
に対応するものと考える。以下も同様である。

\medskip\noindent
最後に、代数空間の二つの写像 $f,g:X\to Y$ が与えられ、これらが $B$ 上にあるとする。
誘導写像 $f,g:X(T)\to Y(T)$ がすべてのスキーム $T$ で $B$ 上にあるものに対して
等しいならば、$f=g$ である。したがって $f,g:X(Z)\to Y(Z)$ も等しい。これは任意の
第3の代数空間 $Z$ で $B$ 上にあるものに対して成り立つ。したがって例えば、群代数空間 $G$ の公理を
$B$ 上で確認するには、$T$-値点上で図式の可換性を確認すれば十分である。ここで $T$ は
$B$ 上のスキームで
ある。以下の Definition \ref{spaces-groupoids-definition-group-space} ではこの方法を用いる。
% P09-ERR-1277
% P09-ERR-1278




\section{同値関係}
\label{spaces-groupoids-section-equivalence-relations}

\noindent
記法については Groupoids, Section
\ref{groupoids-section-equivalence-relations} を参照せよ。

\begin{definition}
\label{spaces-groupoids-definition-equivalence-relation}
% P09-ERR-1279
Section \ref{spaces-groupoids-section-notation} のように $B\to S$ とする。$U$ を $B$ 上の代数空間とする。
\begin{enumerate}
\item $U$ 上の{\it 前関係}を $B$ 上で考えるとは、任意の射
$j:R\to U\times_B U$ で、$B$ 上の代数空間の射であるものをいう。この場合
$t=\text{pr}_0\circ j$ および $s=\text{pr}_1\circ j$ とおくので、
$j=(t,s)$ である。
\item $U$ 上の{\it 関係}を $B$ 上で考えるとは、単射
$j:R\to U\times_B U$ で、$B$ 上の代数空間の射であるものをいう。
\item {\it 前同値関係}とは、前関係 $j:R\to U\times_B U$ で、写像
$j:R(T)\to U(T)\times U(T)$ の像が同値関係となるものをいう。この条件は、すべての
スキーム $T$ で $B$ 上にあるものに対して課す。
\item 代数空間の射 $R\to U\times_B U$ が $B$ 上にあるとする。これを
{\it $U$ 上の同値関係}と $B$ 上でよぶのは、任意の $T$ で $B$ 上にあるものに対して、
その $T$-値点を考えたとき、$R$ の $T$-値点が $U$ の点集合上に同値関係を定める場合、
かつその場合に限る。
\end{enumerate}
\end{definition}

\noindent
言い換えると、同値関係とは、$j$ が関係でもあるような前同値関係である。

\begin{lemma}
\label{spaces-groupoids-lemma-restrict-relation}
% P09-ERR-1280
Section \ref{spaces-groupoids-section-notation} のように $B\to S$ とする。$U$ を $B$ 上の代数空間とする。
$j:R\to U\times_B U$ を前関係とし、$g:U'\to U$ を $B$ 上の代数空間の射とする。
最後に
$$
R'=(U'\times_B U')\times_{U\times_B U}R
\xrightarrow{j'}
U'\times_B U'
$$
とおく。このとき $j'$ は $U'$ 上の前関係であり、$B$ 上にある。$j$ が関係ならば $j'$ も
関係である。$j$ が前同値関係ならば $j'$ も前同値関係である。$j$ が同値関係ならば
$j'$ も同値関係である。
\end{lemma}

\begin{proof}
省略する。
\end{proof}






\begin{definition}
\label{spaces-groupoids-definition-restrict-relation}
% P09-ERR-1281
Section \ref{spaces-groupoids-section-notation} のように $B\to S$ とする。$U$ を $B$ 上の代数空間とする。
$j:R\to U\times_B U$ を前関係とし、$g:U'\to U$ を $B$ 上の代数空間の射とする。
Lemma \ref{spaces-groupoids-lemma-restrict-relation} の前関係 $j':R'\to U'\times_B U'$ を、
前関係 $j$ の $U'$ への{\it 制限}または{\it 引き戻し}という。この状況では
$R'=R|_{U'}$ と書くことがある。
\end{definition}

\begin{lemma}
\label{spaces-groupoids-lemma-pre-equivalence-equivalence-relation-points}
% P09-ERR-1282
Section \ref{spaces-groupoids-section-notation} のように $B\to S$ とする。
$j:R\to U\times_B U$ を $B$ 上の代数空間の前関係とする。$|U|$ 上の関係を規則
$$
x\sim y
\Leftrightarrow
\exists\ r\in|R|:
t(r)=x,
s(r)=y.
$$
によって定める。$j$ が前同値関係ならば、これは同値関係である。
\end{lemma}

\begin{proof}
$x\sim y$ かつ $y\sim z$ と仮定する。$r\in|R|$ を選び、$t(r)=x$、$s(r)=y$ とする。
さらに $r'\in|R|$ を選び、$t(r')=y$、$s(r')=z$ とする。
体 $K$ を、$r$ と $r'$ が射 $r,r':\Spec(K)\to R$ によって表され、
$s\circ r=t\circ r'$ となるように選べる。$x=t\circ r$、
$y=s\circ r=t\circ r'$、および $z=s\circ r'$ と記す。このとき
$x,y,z:\Spec(K)\to U$ である。構成により $(x,y)\in j(R(K))$ かつ
$(y,z)\in j(R(K))$ である。$j$ は前同値関係なので
$(x,z)\in j(R(K))$ でもある。これは明らかに $x\sim z$ を含意する。

\medskip\noindent
$\sim$ が反射的かつ対称的であることの証明は省略する。
\end{proof}















\section{群代数空間}
\label{spaces-groupoids-section-group-spaces}

\noindent
記法については Groupoids, Section
\ref{groupoids-section-group-schemes} を参照せよ。

\begin{definition}
\label{spaces-groupoids-definition-group-space}
% P09-ERR-1283
Section \ref{spaces-groupoids-section-notation} のように $B\to S$ とする。
\begin{enumerate}
\item {\it $B$ 上の群代数空間}とは対 $(G,m)$ であって、$G$ が $B$ 上の代数空間、
$m:G\times_B G\to G$ が $B$ 上の代数空間の射であり、次の性質を満たすものをいう。
すべてのスキーム $T$ で $B$ 上にあるものに対して、対 $(G(T),m)$ は群である。
\item {\it 群代数空間の射 $\psi:(G,m)\to(G',m')$ を $B$ 上で考えたもの}とは、
代数空間の射 $\psi:G\to G'$ で、$B$ 上にあり、すべての $T/B$ に対して誘導写像
$\psi:G(T)\to G'(T)$ が群準同型となるものをいう。
\end{enumerate}
\end{definition}

\noindent
群代数空間 $(G,m)$ を代数空間 $B$ 上で取る。Groupoids, Section
\ref{groupoids-section-group-schemes} の議論により、$B$ 上の代数空間の射として
（単位元）$e:B\to G$ および（逆元）$i:G\to G$ を得る。これらは、すべての $T$ に
対して四つ組 $(G(T),m,e,i)$ が群の公理を満たすような射である。

\medskip\noindent
群代数空間 $(G,m)$、$(G',m')$ を $B$ 上で取る。
$f:G\to G'$ を $B$ 上の代数空間の射とする。定義から、$f$ が
$B$ 上の群代数空間の射であることと、次の図式が可換であることとは同値である。
$$
\xymatrix{
G \times_B G \ar[r]_-{f \times f} \ar[d]_m &
G' \times_B G' \ar[d]^m \\
G \ar[r]^f & G'
}
$$

\begin{lemma}
\label{spaces-groupoids-lemma-base-change-group-space}
% P09-ERR-1284
Section \ref{spaces-groupoids-section-notation} のように $B\to S$ とする。
群代数空間 $(G,m)$ を $B$ 上で取り、代数空間の射 $B'\to B$ を取る。
引戻し $(G_{B'},m_{B'})$ は $B'$ 上の群代数空間である。
\end{lemma}

\begin{proof}
省略する。
\end{proof}



\section{群代数空間の性質}
\label{spaces-groupoids-section-properties-group-spaces}

\noindent
この節では、任意の基礎上で成り立つ群代数空間のいくつかの簡単な性質をまとめる。

\begin{lemma}
\label{spaces-groupoids-lemma-group-scheme-separated}
$S$ をスキームとし、$B$ を $S$ 上の代数空間とする。
$G$ を $B$ 上の群代数空間とする。このとき、$G\to B$ が分離的
（それぞれ準分離的、局所分離的）であることと、単位射 $e:B\to G$ が閉埋入
（それぞれ準コンパクト、埋入）であることとは同値である。
\end{lemma}

\begin{proof}
Morphisms of Spaces, Lemma \ref{spaces-morphisms-lemma-section-immersion}
により、$e$ は、$G\to B$ が分離的（それぞれ準分離的、局所分離的）ならば閉埋入
（それぞれ準コンパクト、埋入）であることを思い出そう。
逆を示すため、次の図式を考える。
$$
\xymatrix{
G \ar[r]_-{\Delta_{G/B}} \ar[d] &
G \times_B G \ar[d]^{(g, g') \mapsto m(i(g), g')} \\
B \ar[r]^e & G
}
$$
代数幾何における関手的観点を用いれば、この図式がカルテシアンであることを示すのは
演習である。言い換えると、$\Delta_{G/B}$ は $e$ の基底変換である。したがって、$e$ が
閉埋入（それぞれ準コンパクト、埋入）ならば、$\Delta_{G/B}$ もそうである。これには
Spaces, Lemma \ref{spaces-lemma-base-change-immersions}
（それぞれ Morphisms of Spaces, Lemma
\ref{spaces-morphisms-lemma-base-change-quasi-compact}、
それぞれ Spaces, Lemma \ref{spaces-lemma-base-change-immersions}）を参照せよ。
\end{proof}

\begin{lemma}
\label{spaces-groupoids-lemma-group-scheme-unramified-or-lqf}
$S$ をスキーム、$B$ を $S$ 上の代数空間、$G$ を $B$ 上の群代数空間とする。
$G\to B$ は局所有限型であると仮定する。このとき、$G\to B$ が非分岐
（それぞれ局所準有限）であることと、$G\to B$ が $e(b)$ において非分岐
（それぞれ準有限）であることがすべての $b\in |B|$ に対して成り立つこととは同値である。
\end{lemma}

\begin{proof}
Morphisms of Spaces, Lemma \ref{spaces-morphisms-lemma-where-unramified}
（それぞれ Morphisms of Spaces, Lemma
\ref{spaces-morphisms-lemma-base-change-quasi-finite-locus}）により、$U\subset G$ であって
$U\to B$ が非分岐（それぞれ局所準有限）となる最大の開部分空間が存在し、$U$ の構成は
基底変換と可換である。したがって、$B=\Spec(k)$ が体のスペクトルである場合に帰着する。
$g\in G(K)$ を拡大 $K/k$ に値を取る点とする。すると、$g$ が $U$ に属するかどうかを
調べるには、$K$ へ基底変換してよい。ゆえに、
$$
G \to \Spec(k)\text{ is unramified at }e
\Leftrightarrow
G \to \Spec(k)\text{ is unramified at }g
$$
を $k$-有理点 $g$ に対して示せば十分である（$g$ と $e$ における準有限性についても
同様である）。$g$ による平行移動は $G$ の $k$ 上の自己同型なので、これは明らかである。
\end{proof}

\begin{lemma}
\label{spaces-groupoids-lemma-open-over-which-unramified-or-lqf}
$S$ をスキーム、$B$ を $S$ 上の代数空間、$G$ を $B$ 上の群代数空間とする。
$G\to B$ は局所有限型であると仮定する。
\begin{enumerate}
\item 最大の開部分空間 $U\subset B$ であって $G_U\to U$ が非分岐となるものが存在し、$U$ の構成は
基底変換と可換である。
\item 最大の開部分空間 $U\subset B$ であって $G_U\to U$ が局所準有限となるものが存在し、$U$ の構成は
基底変換と可換である。
\end{enumerate}
\end{lemma}

\begin{proof}
Morphisms of Spaces, Lemma \ref{spaces-morphisms-lemma-where-unramified}
（それぞれ Morphisms of Spaces, Lemma
\ref{spaces-morphisms-lemma-base-change-quasi-finite-locus}）により、$W\subset G$ であって
$W\to B$ が非分岐（それぞれ局所準有限）となる最大の開部分空間が存在する。さらに、$W$ の
構成は基底変換と可換である。Lemma \ref{spaces-groupoids-lemma-group-scheme-unramified-or-lqf}
により、いずれの場合にも $U=e^{-1}(W)$ である。
\end{proof}













\section{群代数空間の例}
\label{spaces-groupoids-section-examples-group-spaces}

\noindent
% P09-ERR-1285
$G\to S$ が基礎スキーム $S$ 上の群スキームならば、その基底変換 $G_B$ は、任意の代数空間
$B$ で $S$ 上にあるものに対して、Lemma \ref{spaces-groupoids-lemma-base-change-group-space} により
$B$ 上の群代数空間である。
以下の例ではこの事実をたびたび用いる。

\begin{example}[乗法群代数空間]
\label{spaces-groupoids-example-multiplicative-group}
% P09-ERR-1286
Section \ref{spaces-groupoids-section-notation} のように $B\to S$ とする。任意のスキーム $T$ で $B$ 上にあるものに、
構造層の大域切断の単元からなる群 $\Gamma(T,\mathcal{O}_T^*)$ を対応させる関手を考える。
これは群代数空間
$$
\mathbf{G}_{m, B} = B \times_S \mathbf{G}_{m, S}
$$
によって $B$ 上で表現可能である。ここで $\mathbf{G}_{m,S}$ は $S$ 上の乗法群スキームである。
Groupoids, Example \ref{groupoids-example-multiplicative-group} を参照せよ。
\end{example}

\begin{example}[群代数空間としての1の冪根]
\label{spaces-groupoids-example-roots-of-unity}
% P09-ERR-1287
Section \ref{spaces-groupoids-section-notation} のように $B\to S$ とする。$n\in\mathbf{N}$ とする。
任意のスキーム $T$ で $B$ 上にあるものに、$\Gamma(T,\mathcal{O}_T^*)$ の部分群で
$n$ 乗して1となる元からなるものを対応させる関手を考える。
これは群代数空間
$$
\mu_{n, B} = B \times_S \mu_{n, S}
$$
によって $B$ 上で表現可能である。ここで $\mu_{n,S}$ は $n$ 乗根の、$S$ 上の群スキームである。
Groupoids, Example \ref{groupoids-example-roots-of-unity} を参照せよ。
\end{example}

\begin{example}[加法群代数空間]
\label{spaces-groupoids-example-additive-group}
% P09-ERR-1288
Section \ref{spaces-groupoids-section-notation} のように $B\to S$ とする。任意のスキーム $T$ で $B$ 上にあるものに、
構造層の大域切断からなる群 $\Gamma(T,\mathcal{O}_T)$ を対応させる関手を考える。
これは群代数空間
$$
\mathbf{G}_{a, B} = B \times_S \mathbf{G}_{a, S}
$$
によって $B$ 上で表現可能である。ここで $\mathbf{G}_{a,S}$ は $S$ 上の加法群スキームである。
Groupoids, Example \ref{groupoids-example-additive-group} を参照せよ。
\end{example}

\begin{example}[一般線形群代数空間]
\label{spaces-groupoids-example-general-linear-group}
% P09-ERR-1289
Section \ref{spaces-groupoids-section-notation} のように $B\to S$ とする。$n\geq 1$ とする。
任意のスキーム $T$ で $B$ 上にあるものに、群
$$
\text{GL}_n(\Gamma(T, \mathcal{O}_T))
$$
すなわち構造層の大域切断上の可逆 $n\times n$ 行列からなる群を対応させる関手を考える。
これは群代数空間
$$
\text{GL}_{n, B} = B \times_S \text{GL}_{n, S}
$$
によって $B$ 上で表現可能である。
% P09-ERR-1290
ここで $\mathbf{G}_{m,S}$ は $S$ 上の一般線形群スキームである。
Groupoids, Example \ref{groupoids-example-general-linear-group} を参照せよ。
\end{example}

\begin{example}
\label{spaces-groupoids-example-determinant}
% P09-ERR-1291
Section \ref{spaces-groupoids-section-notation} のように $B\to S$ とする。$n\geq 1$ とする。
行列式は群代数空間の射
$$
\det : \text{GL}_{n, B} \longrightarrow \mathbf{G}_{m, B}
$$
を $B$ 上で定める。これは $S$ 上の行列式射、すなわち
Groupoids, Example \ref{groupoids-example-determinant} の射の基底変換である。
\end{example}

\begin{example}[定数群代数空間]
\label{spaces-groupoids-example-constant-group}
% P09-ERR-1292
Section \ref{spaces-groupoids-section-notation} のように $B\to S$ とする。$G$ を抽象群とする。
任意のスキーム $T$ で $B$ 上にあるものに、局所定数写像 $T\to G$ の群を対応させる関手を考える
（ここで $T$ には Zariski 位相を、$G$ には離散位相を入れる）。これは群代数空間
$$
G_B = B \times_S G_S
$$
によって $B$ 上で表現可能である。ここで $G_S$ は Groupoids, Example
\ref{groupoids-example-constant-group} で導入した定数群スキームである。
\end{example}






\section{群代数空間の作用}
\label{spaces-groupoids-section-action-group-space}

\noindent
記法については Groupoids, Section \ref{groupoids-section-action-group-scheme} を参照せよ。

\begin{definition}
\label{spaces-groupoids-definition-action-group-space}
% P09-ERR-1293
Section \ref{spaces-groupoids-section-notation} のように $B\to S$ とする。
群代数空間 $(G,m)$ を $B$ 上で取り、代数空間 $X$ を $B$ 上で取る。
\begin{enumerate}
\item {\it $G$ の代数空間 $X/B$ への作用}とは、射 $a:G\times_B X\to X$ であって
$B$ 上にあり、すべてのスキーム $T$ で $B$ 上にあるものに対し、写像
$a:G(T)\times X(T)\to X(T)$ が $G(T)$-集合としての構造を $X(T)$ に定めるものをいう。
\item $X$、$Y$ を $B$ 上の代数空間とし、それぞれ $G$ の作用を備えるものとする。
{\it 同変}、より正確には {\it $G$-同変}な射 $\psi:X\to Y$ とは、$B$ 上の代数空間の射であって、
すべての $T$ で $B$ 上にあるものに対し、写像
$\psi:X(T)\to Y(T)$ が $G(T)$-集合の射であることをいう。
\end{enumerate}
\end{definition}

\noindent
(1) の状況では、これは図式
\begin{equation}
\label{spaces-groupoids-equation-action}
\xymatrix{
G \times_B G \times_B X \ar[r]_-{1_G \times a} \ar[d]_{m \times 1_X} &
G \times_B X \ar[d]^a \\
G \times_B X \ar[r]^a & X
}
\quad
\xymatrix{
G \times_B X \ar[r]_-a & X \\
X\ar[u]^{e \times 1_X} \ar[ru]_{1_X}
}
\end{equation}
が可換であることを意味する。(2) の状況では、単に次の図式が可換であることを意味する。
% P09-ERR-1294
$$
\xymatrix{
G \times_B X \ar[r]_-{\text{id} \times f} \ar[d]_a &
G \times_B Y \ar[d]^a \\
X \ar[r]^f & Y
}
$$

\begin{definition}
\label{spaces-groupoids-definition-free-action}
$B\to S$、$G\to B$、$X\to B$ を Definition \ref{spaces-groupoids-definition-action-group-space}
のように取る。$a:G\times_B X\to X$ を $G$ の $X/B$ への作用とする。
すべてのスキーム $T$ で $B$ 上にあるものに対して、作用
$a:G(T)\times X(T)\to X(T)$ が群 $G(T)$ の集合 $X(T)$ への自由作用であるとき、
この作用を {\it 自由}という。
\end{definition}

\begin{lemma}
\label{spaces-groupoids-lemma-free-action}
% P09-ERR-1295
Definition \ref{spaces-groupoids-definition-free-action} の状況とする。作用 $a$ が自由であることと、
$$
G \times_B X \to X \times_B X, \quad (g, x) \mapsto (a(g, x), x)
$$
が代数空間のモノ射であることとは同値である。
\end{lemma}

\begin{proof}
定義から直ちに従う。
\end{proof}





\section{主等質空間}
\label{spaces-groupoids-section-principal-homogeneous}

\noindent
この節は Groupoids, Section \ref{groupoids-section-principal-homogeneous} の類似である。
先にそちらを読むことを勧める。

\begin{definition}
\label{spaces-groupoids-definition-pseudo-torsor}
$S$ をスキームとし、$B$ を $S$ 上の代数空間とする。
群代数空間 $(G,m)$ を $B$ 上で取る。代数空間 $X$ を $B$ 上で取り、
$a:G\times_B X\to X$ を $G$ の $X$ への作用とする。
\begin{enumerate}
\item $X$ が {\it 擬 $G$-トーサー}である、または $X$ が
{\it $G$ の下で形式的に主等質}であるとは、誘導される射
$G\times_B X\to X\times_B X$、$(g,x)\mapsto(a(g,x),x)$ が同型であることをいう。
\item % P09-ERR-1296
擬 $G$-トーサー $X$ が {\it 自明}であるとは、$G$-同変同型 $G\to X$ が
$B$ 上に存在することをいう。ここで $G$ は $G$ に左乗法で作用する。
\end{enumerate}
\end{definition}

\noindent
代数空間の射 $B'\to B$ があるとき、引戻し $X_{B'}$ は、擬 $G$-トーサーで
$B$ 上にあるものから得られるならば、擬 $G_{B'}$-トーサーで $B'$ 上にあることは明らかである。

\begin{lemma}
\label{spaces-groupoids-lemma-characterize-trivial-pseudo-torsors}
Definition \ref{spaces-groupoids-definition-pseudo-torsor} の状況とする。
\begin{enumerate}
\item 代数空間 $X$ が擬 $G$-トーサーであることと、すべてのスキーム $T$ で $B$ 上に
あるものに対し、集合 $X(T)$ が空であるか、または群 $G(T)$ の $X(T)$ への作用が
単純推移的であることとは同値である。
\item 擬 $G$-トーサー $X$ が自明であることと、射 $X\to B$ が切断をもつこととは同値である。
\end{enumerate}
\end{lemma}

\begin{proof}
省略する。
\end{proof}

\begin{definition}
\label{spaces-groupoids-definition-principal-homogeneous-space}
$S$ をスキームとする。$B$ を $S$ 上の代数空間とする。
群代数空間 $(G,m)$ を $B$ 上で取る。$X$ を擬 $G$-トーサーで $B$ 上にあるものとする。
\begin{enumerate}
\item $X$ が {\it 主等質空間}、より正確には
{\it 主等質 $G$-空間で $B$ 上にあるもの}であるとは、fpqc 被覆\footnote{Groupoids, Definition
\ref{groupoids-definition-principal-homogeneous-space} における既定のトーサーは、
fpqc 被覆上で自明となる擬トーサーである。
% P09-ERR-1297
$G$ は代数空間として群の層とみなせるので、すでに $G$-トーサーという概念があり、
これは fppf トーサーに対応する。Lemma \ref{spaces-groupoids-lemma-torsor} を参照せよ。
したがって、fpqc 局所的に自明な擬トーサーには「主等質空間」という語を用い、
この状況ではトーサーという語を避けるようにする。}
$\{B_i\to B\}_{i\in I}$ であって、各 $X_{B_i}\to B_i$ が切断をもつ
（すなわち自明な擬 $G_{B_i}$-トーサーである）ものが存在することをいう。
\item $\tau\in\{Zariski,\etale,smooth,syntomic,fppf\}$ とする。
$X$ が {\it $G$-トーサーで $\tau$ 位相におけるもの}、または {\it $\tau$ $G$-トーサー}、
あるいは単に {\it $\tau$ トーサー}であるとは、$\tau$ 被覆
$\{B_i\to B\}_{i\in I}$ であって、各 $X_{B_i}\to B_i$ が切断をもつものが存在することをいう。
\item $X$ が主等質 $G$-空間で $B$ 上にあるとき、それが \'etale 位相のトーサーならば
{\it 準等自明}であるという。
\item $X$ が主等質 $G$-空間で $B$ 上にあるとき、それが Zariski 位相のトーサーならば
{\it 局所自明}であるという。
\end{enumerate}
\end{definition}

\noindent
「$X$ を $G$-主等質空間で $B$ 上にあるものとする」と言うことがある。これは、$X$ が
$B$ 上の代数空間で、$G$ の作用を備え、それによって $B$ 上の主等質空間となることを表す。
次に、両方が適用できる場合には、この用語法が先に導入したものと一致することを示す。

\begin{lemma}
\label{spaces-groupoids-lemma-torsor}
$S$ をスキームとする。群代数空間 $(G,m)$ を $S$ 上で取る。
代数空間 $X$ を $S$ 上で取り、$a:G\times_S X\to X$ を $G$ の $X$ への作用とする。
このとき、$X$ が $G$-トーサーで $fppf$-位相におけるものとなることを Definition
\ref{spaces-groupoids-definition-principal-homogeneous-space} の意味で考えたものと、$X$ が $G$-トーサーで
$(\Sch/S)_{fppf}$ 上にあるものとなることを Cohomology on Sites, Definition
\ref{sites-cohomology-definition-torsor} の意味で考えたものとは同値である。
\end{lemma}

\begin{proof}
省略する。
\end{proof}

\begin{lemma}
\label{spaces-groupoids-lemma-pseudo-torsor-implications}
$S$ をスキームとし、$B$ を $S$ 上の代数空間とする。
群代数空間 $G$ を $B$ 上で取る。$X$ を擬 $G$-トーサーで $B$ 上にあるものとする。
% P09-ERR-1298
$G$ および $X$ は $B$ 上局所有限型であると仮定する。
\begin{enumerate}
\item $G\to B$ が非分岐ならば、$X\to B$ は非分岐である。
\item $G\to B$ が局所準有限ならば、$X\to B$ は局所準有限である。
\end{enumerate}
\end{lemma}

\begin{proof}
(1) を証明する。Morphisms of Spaces, Lemma
\ref{spaces-morphisms-lemma-where-unramified} により、$B$ が体のスペクトルである場合に
帰着する。$X$ が空ならば結論は成り立つ。$X$ が空でなければ、体を拡大した後、
$X$ が点をもつと仮定してよい。このとき $G\cong X$ であり、結論が従う。

\medskip\noindent
(2) の証明も Morphisms of Spaces, Lemma
\ref{spaces-morphisms-lemma-base-change-quasi-finite-locus} を用いればまったく同様に進む。
\end{proof}
















\section{同変準連接層}
\label{spaces-groupoids-section-equivariant}

\noindent
Groupoids, Section \ref{groupoids-section-equivariant} と比較せよ。

\begin{definition}
\label{spaces-groupoids-definition-equivariant-module}
Section \ref{spaces-groupoids-section-notation} のように $B\to S$ とする。群代数空間 $(G,m)$ を $B$ 上で取り、
$a:G\times_B X\to X$ を $G$ の代数空間 $X$ への作用とし、この空間は $B$ 上にあるものとする。
{\it $G$-同変準連接 $\mathcal{O}_X$-加群}、または単に
{\it 同変準連接 $\mathcal{O}_X$-加群}とは、対 $(\mathcal{F},\alpha)$ であって、
$\mathcal{F}$ が準連接 $\mathcal{O}_X$-加群、$\alpha$ が $\mathcal{O}_{G\times_B X}$-加群写像
$$
\alpha : a^*\mathcal{F} \longrightarrow \text{pr}_1^*\mathcal{F}
$$
であるものをいう。ここで $\text{pr}_1:G\times_B X\to X$ は射影であり、次を満たすものとする。
\begin{enumerate}
\item 図式
$$
\xymatrix{
(1_G \times a)^*\text{pr}_2^*\mathcal{F} \ar[r]_-{\text{pr}_{12}^*\alpha} &
\text{pr}_2^*\mathcal{F} \\
(1_G \times a)^*a^*\mathcal{F} \ar[u]^{(1_G \times a)^*\alpha} \ar@{=}[r] &
(m \times 1_X)^*a^*\mathcal{F} \ar[u]_{(m \times 1_X)^*\alpha}
}
$$
% P09-ERR-1299
が $\mathcal{O}_{G\times_B G\times_B X}$-加群の圏で可換である。
\item 引戻し
$$
(e \times 1_X)^*\alpha : \mathcal{F} \longrightarrow \mathcal{F}
$$
が恒等写像である。
\end{enumerate}
説明については Equation (\ref{spaces-groupoids-equation-action}) の対応する図式と比較せよ。
\end{definition}

\noindent
第1の図式の可換性により、$(e\times 1_X)^*\alpha$ は $\mathcal{F}$ 上の冪等作用素となる。
したがって、条件 (2) は単にそれが同型であるという条件であることに注意せよ。

\begin{lemma}
\label{spaces-groupoids-lemma-pullback-equivariant}
Section \ref{spaces-groupoids-section-notation} のように $B\to S$ とする。
群代数空間 $G$ を $B$ 上で取る。$f:X\to Y$ を $G$-同変射であって、
$B$ 上の代数空間の間にあり、それらが $G$-作用を備えるものとする。
このとき、引戻し $f^*$ であって
$(\mathcal{F},\alpha)\mapsto(f^*\mathcal{F},(1_G\times f)^*\alpha)$
で与えられるものは、準連接 $G$-同変層で $Y$ 上にあるものの圏から
準連接 $G$-同変層で $X$ 上にあるものの圏への関手を定める。
\end{lemma}

\begin{proof}
省略する。
\end{proof}






\section{代数空間における亜群}
\label{spaces-groupoids-section-groupoids}

\noindent
記法については Groupoids, Section \ref{groupoids-section-groupoids} を参照せよ。

\begin{definition}
\label{spaces-groupoids-definition-groupoid}
Section \ref{spaces-groupoids-section-notation} のように $B\to S$ とする。
\begin{enumerate}
\item {\it $B$ 上の代数空間における亜群}とは五つ組 $(U,R,s,t,c)$ であって、
$U$ および $R$ が $B$ 上の代数空間であり、
$s,t:R\to U$ および $c:R\times_{s,U,t}R\to R$ が $B$ 上の代数空間の射で、
次の性質を満たすものをいう。任意のスキーム $T$ で $B$ 上にあるものに対して、五つ組
$$
(U(T), R(T), s, t, c)
$$
は亜群圏である。
\item {\it 代数空間における亜群の射
$f:(U,R,s,t,c)\to(U',R',s',t',c')$ を $B$ 上で考えたもの}とは、
代数空間の射 $f:U\to U'$ および $f:R\to R'$ で $B$ 上にあるものによって与えられ、
次の性質を満たすものをいう。任意のスキーム $T$ で $B$ 上にあるものに対し、写像 $f$ は
亜群圏 $(U(T),R(T),s,t,c)$ から亜群圏 $(U'(T),R'(T),s',t',c')$ への関手を定める。
\end{enumerate}
\end{definition}

\noindent
$(U,R,s,t,c)$ を $B$ 上の代数空間における亜群とする。代数空間の射
$e:U\to R$ および $i:R\to R$ で $B$ 上にあり、すべてのスキーム $T$ で $B$ 上にあるものに
対し、誘導写像 $e:U(T)\to R(T)$ が恒等射を、$i:R(T)\to R(T)$ が亜群圏における逆射を
与えるものが一意に存在することに注意せよ。七つ組 $(U,R,s,t,c,e,i)$ は、
Groupoids, Section \ref{groupoids-section-groupoids} の公理 (1), (2)(a), (2)(b), (3)(a),
(3)(b) のそれぞれに対応する可換図式を満たす。逆に、この性質をもつ七つ組が与えられれば、
五つ組 $(U,R,s,t,c)$ は $B$ 上の代数空間における亜群である。$i$ は同型であり、$e$ は
$s$ と $t$ のいずれの切断でもあることに注意せよ。さらに、$B$ 上の代数空間における亜群が
与えられたとき、
$$
j = (t, s) : R \longrightarrow U \times_B U
$$
と記す。これは上の Section \ref{spaces-groupoids-section-equivalence-relations} における規約と整合する。
単位射と逆射の存在を強調するため、「$(U,R,s,t,c,e,i)$ を $B$ 上の代数空間における亜群とする」
と言うことがある。

\begin{lemma}
\label{spaces-groupoids-lemma-groupoid-pre-equivalence}
Section \ref{spaces-groupoids-section-notation} のように $B\to S$ とする。
代数空間における亜群 $(U,R,s,t,c)$ を $B$ 上で取ると、射
$j:R\to U\times_B U$ は前同値関係である。
\end{lemma}

\begin{proof}
省略する。これは定義に関するよい演習である。
\end{proof}

\begin{lemma}
\label{spaces-groupoids-lemma-equivalence-groupoid}
Section \ref{spaces-groupoids-section-notation} のように $B\to S$ とする。
同値関係 $j:R\to U\times_B U$ を $B$ 上で取ると、これを代数空間における亜群
$(U,R,s,t,c)$ で $B$ 上にあるものへ拡張する方法が一意に存在する。
\end{lemma}

\begin{proof}
省略する。これは定義に関するよい演習である。
\end{proof}

\begin{lemma}
\label{spaces-groupoids-lemma-diagram}
Section \ref{spaces-groupoids-section-notation} のように $B\to S$ とする。
$(U,R,s,t,c)$ を $B$ 上の代数空間における亜群とする。可換図式
$$
\xymatrix{
& U & \\
R \ar[d]_s \ar[ru]^t &
R \times_{s, U, t} R
\ar[l]^-{\text{pr}_0} \ar[d]^{\text{pr}_1} \ar[r]_-c &
R \ar[d]^s \ar[lu]_t \\
U & R \ar[l]_t \ar[r]^s & U
}
$$
において、下側の二つの正方形はファイバー積図式である。さらに、上側の三角形
（実際には正方形である）もカルテシアンである。
\end{lemma}

\begin{proof}
省略する。定義と代数幾何における関手的観点に関する演習である。
\end{proof}

\begin{lemma}
\label{spaces-groupoids-lemma-diagram-pull}
Section \ref{spaces-groupoids-section-notation} のように $B\to S$ とする。
$(U,R,s,t,c,e,i)$ を $B$ 上の代数空間における亜群とする。図式
\begin{equation}
\label{spaces-groupoids-equation-pull}
\xymatrix{
R \times_{t, U, t} R
\ar@<1ex>[r]^-{\text{pr}_1} \ar@<-1ex>[r]_-{\text{pr}_0}
\ar[d]_{\text{pr}_0 \times c \circ (i, 1)} &
R \ar[r]^t \ar[d]^{\text{id}_R} &
U \ar[d]^{\text{id}_U} \\
R \times_{s, U, t} R
\ar@<1ex>[r]^-c \ar@<-1ex>[r]_-{\text{pr}_0} \ar[d]_{\text{pr}_1} &
R \ar[r]^t \ar[d]^s &
U \\
R \ar@<1ex>[r]^s \ar@<-1ex>[r]_t &
U
}
\end{equation}
は可換である。上側の二つの行は、表示された鉛直写像によって同型である。
左下の二つの正方形はカルテシアンである。
\end{lemma}

\begin{proof}
図式の可換性は亜群の公理から従う。亜群の言葉では、左上の鉛直矢印は、終域が等しい
射の対 $(\alpha,\beta)$ に、射の対 $(\alpha,\alpha^{-1}\circ\beta)$ を対応させる。
任意の亜群において、これは
$\text{Arrows}\times_{t,\text{Ob},t}\text{Arrows}$ と
$\text{Arrows}\times_{s,\text{Ob},t}\text{Arrows}$ の間の全単射を定める。
したがって補題の第2の主張が従う。最後の主張は Lemma \ref{spaces-groupoids-lemma-diagram} から従う。
\end{proof}

\begin{lemma}
\label{spaces-groupoids-lemma-base-change-groupoid}
Section \ref{spaces-groupoids-section-notation} のように $B\to S$ とする。
$(U,R,s,t,c)$ を $B$ 上の代数空間における亜群とする。
代数空間の射 $B'\to B$ を取る。このとき、基底変換
$U'=B'\times_B U$、$R'=B'\times_B R$ に、基底変換 $s'$、$t'$、$c'$ であって
射 $s,t,c$ のものを備えたものは、代数空間における亜群
$(U',R',s',t',c')$ で $B'$ 上にあるものをなし、射影は代数空間における亜群の射
$$
(U', R', s', t', c') \to (U, R, s, t, c)
$$
で $B$ 上にあるものを定める。
\end{lemma}

\begin{proof}
省略する。ヒント：
$R'\times_{s',U',t'}R'=B'\times_B(R\times_{s,U,t}R)$。
\end{proof}





\section{亜群上の準連接層}
\label{spaces-groupoids-section-groupoids-quasi-coherent}

\noindent
Groupoids, Section \ref{groupoids-section-groupoids-quasi-coherent} と比較せよ。

\begin{definition}
\label{spaces-groupoids-definition-groupoid-module}
Section \ref{spaces-groupoids-section-notation} のように $B\to S$ とする。
$(U,R,s,t,c)$ を $B$ 上の代数空間における亜群とする。
{\it $(U,R,s,t,c)$ 上の準連接加群}とは対 $(\mathcal{F},\alpha)$ であって、
$\mathcal{F}$ が準連接 $\mathcal{O}_U$-加群、$\alpha$ が $\mathcal{O}_R$-加群写像
$$
\alpha : t^*\mathcal{F} \longrightarrow s^*\mathcal{F}
$$
であり、次を満たすものをいう。
\begin{enumerate}
\item 図式
$$
\xymatrix{
& \text{pr}_1^*t^*\mathcal{F} \ar[r]_-{\text{pr}_1^*\alpha} &
\text{pr}_1^*s^*\mathcal{F} \ar@{=}[rd] & \\
\text{pr}_0^*s^*\mathcal{F} \ar@{=}[ru] & & & c^*s^*\mathcal{F} \\
& \text{pr}_0^*t^*\mathcal{F} \ar[lu]^{\text{pr}_0^*\alpha} \ar@{=}[r] &
c^*t^*\mathcal{F} \ar[ru]_{c^*\alpha}
}
$$
% P09-ERR-1300
が $\mathcal{O}_{R\times_{s,U,t}R}$-加群の圏で可換である。
\item 引戻し
$$
e^*\alpha : \mathcal{F} \longrightarrow \mathcal{F}
$$
が恒等写像である。
\end{enumerate}
Lemma \ref{spaces-groupoids-lemma-diagram} の可換図式と比較せよ。
\end{definition}

\noindent
第1の図式の可換性により、作用素 $e^*\alpha$ は冪等となる。したがって、第2の条件は
$e^*\alpha$ が同型であるという条件に言い換えられる。実際、この条件から $\alpha$ は同型となる。

\begin{lemma}
\label{spaces-groupoids-lemma-isomorphism}
Section \ref{spaces-groupoids-section-notation} のように $B\to S$ とする。
$(U,R,s,t,c)$ を $B$ 上の代数空間における亜群とする。
$(\mathcal{F},\alpha)$ が $(U,R,s,t,c)$ 上の準連接加群ならば、$\alpha$ は同型である。
\end{lemma}

\begin{proof}
Definition \ref{spaces-groupoids-definition-groupoid-module} の可換図式を射
$(i,1):R\to R\times_{s,U,t}R$ で引き戻す。このとき
$i^*\alpha\circ\alpha=s^*e^*\alpha$ であることがわかる。
射 $(1,i)$ で引き戻すと、関係
$\alpha\circ i^*\alpha=t^*e^*\alpha$ を得る。第2の仮定により、これらの射は恒等射である。
したがって $i^*\alpha$ は $\alpha$ の逆射である。
\end{proof}

\begin{lemma}
\label{spaces-groupoids-lemma-pullback}
Section \ref{spaces-groupoids-section-notation} のように $B\to S$ とする。射
$$
f : (U, R, s, t, c) \to (U', R', s', t', c')
$$
% P09-ERR-1301
を、代数空間における亜群の射で $B$ 上にあるものとする。このとき、引戻し $f^*$ であって
$$
(\mathcal{F}, \alpha) \mapsto (f^*\mathcal{F}, f^*\alpha)
$$
で与えられるものは、$(U',R',s',t',c')$ 上の準連接層の圏から
$(U,R,s,t,c)$ 上の準連接層の圏への関手を定める。
\end{lemma}

\begin{proof}
省略する。
\end{proof}

\begin{lemma}
\label{spaces-groupoids-lemma-pushforward}
Section \ref{spaces-groupoids-section-notation} のように $B\to S$ とする。射
$$
f : (U, R, s, t, c) \to (U', R', s', t', c')
$$
を、代数空間における亜群の射で $B$ 上にあるものとする。次を仮定する。
\begin{enumerate}
\item $f:U\to U'$ は準コンパクトかつ準分離的である。
\item 正方形
$$
\xymatrix{
R \ar[d]_t \ar[r]_f & R' \ar[d]^{t'} \\
U \ar[r]^f & U'
}
$$
はカルテシアンである。
\item $s'$ および $t'$ は平坦である。
\end{enumerate}
このとき、順像 $f_*$ であって
$$
(\mathcal{F}, \alpha) \mapsto (f_*\mathcal{F}, f_*\alpha)
$$
で与えられるものは、$(U,R,s,t,c)$ 上の準連接層の圏から
$(U',R',s',t',c')$ 上の準連接層の圏への関手を定め、これは
Lemma \ref{spaces-groupoids-lemma-pullback} で定義した引戻しの右随伴である。
\end{lemma}

\begin{proof}
$U\to U'$ は準コンパクトかつ準分離的なので、$f_*$ は準連接層を準連接層へ移す
（Morphisms of Spaces, Lemma \ref{spaces-morphisms-lemma-pushforward}）。
さらに、正方形
$$
\vcenter{
\xymatrix{
R \ar[d]_t \ar[r]_f & R' \ar[d]^{t'} \\
U \ar[r]^f & U'
}
}
\quad\text{and}\quad
\vcenter{
\xymatrix{
R \ar[d]_s \ar[r]_f & R' \ar[d]^{s'} \\
U \ar[r]^f & U'
}
}
$$
はカルテシアンなので、
$(t')^*f_*\mathcal{F}=f_*t^*\mathcal{F}$ および
$(s')^*f_*\mathcal{F}=f_*s^*\mathcal{F}$ を得る。Cohomology of Spaces, Lemma
\ref{spaces-cohomology-lemma-flat-base-change-cohomology} を参照せよ。
したがって、$f_*\alpha$ を写像
$(t')^*f_*\mathcal{F}\to(s')^*f_*\mathcal{F}$ とみなすことには意味がある。同様の議論により、
$f_*\alpha$ は余サイクル条件を満たす。この関手が引戻し関手の随伴となるのは、
環付き空間上の加群の引戻しと順像が随伴であることによる。いくつかの詳細は省略する。
\end{proof}


\begin{lemma}
\label{spaces-groupoids-lemma-colimits}
Section \ref{spaces-groupoids-section-notation} のように $B\to S$ とする。
$(U,R,s,t,c)$ を $B$ 上の代数空間における亜群とする。
$(U,R,s,t,c)$ 上の準連接加群の圏は余極限をもつ。
\end{lemma}

\begin{proof}
$i\mapsto(\mathcal{F}_i,\alpha_i)$ を添字圏 $\mathcal{I}$ 上の図式とする。余極限
$\mathcal{F}=\colim\mathcal{F}_i$ を作ることができ、これは $U$ 上の準連接層である。
Properties of Spaces, Lemma
\ref{spaces-properties-lemma-properties-quasi-coherent} を参照せよ。
余極限は引戻しと可換なので、$s^*\mathcal{F}=\colim s^*\mathcal{F}_i$ であり、同様に
$t^*\mathcal{F}=\colim t^*\mathcal{F}_i$ である。したがって
$\alpha=\colim\alpha_i$ とおける。$(\mathcal{F},\alpha)$ が
$(U,R,s,t,c)$ 上の準連接加群の圏におけるこの図式の余極限であることの証明は省略する。
\end{proof}

\begin{lemma}
\label{spaces-groupoids-lemma-abelian}
Section \ref{spaces-groupoids-section-notation} のように $B\to S$ とする。
$(U,R,s,t,c)$ を $B$ 上の代数空間における亜群とする。
$s$、$t$ が平坦ならば、$(U,R,s,t,c)$ 上の準連接加群の圏は Abel 圏である。
\end{lemma}

\begin{proof}
$\varphi:(\mathcal{F},\alpha)\to(\mathcal{G},\beta)$ を
$(U,R,s,t,c)$ 上の準連接加群の準同型とする。$s$ は平坦なので、
$$
0 \to s^*\Ker(\varphi)
\to s^*\mathcal{F} \to s^*\mathcal{G} \to s^*\Coker(\varphi) \to 0
$$
は完全であり、$t$ による引戻しについても同様である。したがって $\alpha$ と $\beta$ は、
余サイクル条件を満たす同型
$\kappa:t^*\Ker(\varphi)\to s^*\Ker(\varphi)$ および
$\lambda:t^*\Coker(\varphi)\to s^*\Coker(\varphi)$ を誘導する。このとき、
$(\Ker(\varphi),\kappa)$ と $(\Coker(\varphi),\lambda)$ が
$(U,R,s,t,c)$ 上の準連接加群の圏における核と余核であることは直ちに確認できる。
さらに、$\Coim(\varphi)=\Im(\varphi)$ という条件は $U$ 上で成り立つので従う。
\end{proof}











\section{準連接加群の余極限}
\label{spaces-groupoids-section-colimits}

\noindent
この節は Groupoids, Section \ref{groupoids-section-colimits} の類似である。

\begin{lemma}
\label{spaces-groupoids-lemma-construct-quasi-coherent}
Section \ref{spaces-groupoids-section-notation} のように $B\to S$ とする。
$(U,R,s,t,c)$ を $B$ 上の代数空間における亜群とする。$s,t$ は平坦、準コンパクト、
かつ準分離的であると仮定する。任意の準連接加群 $\mathcal{G}$ で $U$ 上にあるものに対し、標準同型
$\alpha:t^*s_*t^*\mathcal{G}\to s^*s_*t^*\mathcal{G}$ が存在し、これによって
$(s_*t^*\mathcal{G},\alpha)$ は $(U,R,s,t,c)$ 上の準連接加群となる。この構成は関手
$$
\QCoh(\mathcal{O}_U) \longrightarrow \QCoh(U, R, s, t, c)
$$
を定め、これは忘却関手 $(\mathcal{F},\beta)\mapsto\mathcal{F}$ の右随伴である。
\end{lemma}

\begin{proof}
準コンパクトかつ準分離的な射に沿う準連接加群の順像は準連接である。
Morphisms of Spaces, Lemma \ref{spaces-morphisms-lemma-pushforward} を参照せよ。
したがって $s_*t^*\mathcal{G}$ は準連接である。Lemma \ref{spaces-groupoids-lemma-diagram} の記法により、
$$
t^*s_*t^*\mathcal{G} =
\text{pr}_{1, *}\text{pr}_0^*t^*\mathcal{G} =
\text{pr}_{1, *}c^*t^*\mathcal{G} =
s^*s_*t^*\mathcal{G}
$$
を得る。中央の等式は $t\circ c=t\circ\text{pr}_0$ が射
$R\times_{s,U,t}R\to U$ として成り立つためであり、最初と最後の等式は、これらの段階で
基底変換と順像が可換であることが Cohomology of Spaces, Lemma
\ref{spaces-cohomology-lemma-flat-base-change-cohomology} からわかるためである。

\medskip\noindent
Definition \ref{spaces-groupoids-definition-groupoid-module} における $\alpha$ の余サイクル条件と随伴性を確認するため、
構成 $\mathcal{G}\mapsto(s_*t^*\mathcal{G},\alpha)$ を別の方法で記述する。
% P09-ERR-1302
亜群スキーム $(R,R\times_{t,U,t}R,\text{pr}_0,\text{pr}_1,\text{pr}_{02})$ であって、
同値関係 $R\times_{t,U,t}R$ から $R$ 上に得られるものを考える。
Lemma \ref{spaces-groupoids-lemma-equivalence-groupoid} を参照せよ。亜群スキームの射
$$
f :
(R, R \times_{t, U, t} R, \text{pr}_1, \text{pr}_0, \text{pr}_{02})
\longrightarrow
(U, R, s, t, c)
$$
が存在する。これは $s:R\to U$ および
$R\times_{t,U,t}R\to R$ で $(r_0,r_1)\mapsto r_0^{-1}\circ r_1$ によって
与えられるものからなる。必要な図式の可換性の確認は省略する。
$t,s:R\to U$ は準コンパクト、準分離的、かつ平坦であり、さらにカルテシアン正方形
$$
\xymatrix{
R \times_{t, U, t} R \ar[d]_{\text{pr}_0}
\ar[rr]_-{(r_0, r_1) \mapsto r_0^{-1} \circ r_1} & & R \ar[d]^t \\
R \ar[rr]^s & & U
}
$$
があるので、Lemma \ref{spaces-groupoids-lemma-diagram-pull} により Lemma \ref{spaces-groupoids-lemma-pushforward} は $f$ に適用できる。
したがって、$f$ に沿う準連接加群の順像と引戻しは随伴関手である。証明を終えるため、これらの関手を
上で記述した関手と同一視する。そのため、次のことに注意する。
$$
t^* :
\QCoh(\mathcal{O}_U)
\longrightarrow
\QCoh(R, R \times_{t, U, t} R, \text{pr}_1, \text{pr}_0, \text{pr}_{02})
$$
は、$\{t:R\to U\}$ が fpqc 被覆なので、準連接層の降下理論により同値である。
Descent on Spaces, Proposition
\ref{spaces-descent-proposition-fpqc-descent-quasi-coherent} を参照せよ。

\medskip\noindent
$f$ に沿う順像と同値 $t^*$ の合成は、$\mathcal{G}$ を
$(s_*t^*\mathcal{G},\alpha)$ へ送る。この方法で得られる同型 $\alpha$ が上で構成したものと
同じであることの確認は省略する。

\medskip\noindent
$f$ に沿う引戻しと同値 $t^*$ の逆との合成は、$(\mathcal{F},\beta)$ を、
$\{t:R\to U\}$ に関する $s^*\mathcal{F}$ の降下へ送る。ここで
降下データ $\gamma$ で $R\times_{t,U,t}R$ 上にあるものは、$\beta$ を
$(r_0,r_1)\mapsto r_0^{-1}\circ r_1$ で引き戻したものである。
同型 $\beta:t^*\mathcal{F}\to s^*\mathcal{F}$ を考える。
$t^*\mathcal{F}$ 上の標準降下データで $\{t:R\to U\}$ に関するもの
（Descent on Spaces, Definition
\ref{spaces-descent-definition-descent-datum-effective-quasi-coherent}）は、
$\beta$ を介して写像
$$
\text{pr}_0^*s^*\mathcal{F}
\xrightarrow{\text{pr}_0^*\beta^{-1}}
\text{pr}_0^*t^*\mathcal{F}
\xrightarrow{can}
\text{pr}_1^*t^*\mathcal{F}
\xrightarrow{\text{pr}_1^*\beta}
\text{pr}_1^*s^*\mathcal{F}
$$
へ移る。$\beta$ は余サイクル条件を満たすので、これは $\beta$ を
$(r_0,r_1)\mapsto r_0^{-1}\circ r_1$ で引き戻したものに等しい。これを見るには、
Definition \ref{spaces-groupoids-definition-groupoid-module} の実際の余サイクル関係を射
$(\text{pr}_0,c\circ(i,1)):R\times_{t,U,t}R\to R\times_{s,U,t}R$
で引き戻せばよい。この射は Lemma \ref{spaces-groupoids-lemma-diagram-pull} の可換図式でも役割を果たす。
したがって、$(s^*\mathcal{F},\gamma)$ は $(t^*\mathcal{F},can)$ と同型である。
以上により、$f$ による引戻しと同値 $t^*$ の逆との合成は、忘却関手
$(\mathcal{F},\beta)\mapsto\mathcal{F}$ と同型である。
\end{proof}

\begin{remark}
\label{spaces-groupoids-remark-adjunction-map}
Lemma \ref{spaces-groupoids-lemma-construct-quasi-coherent} の状況で、忘却関手を
$$
F : \QCoh(U, R, s, t, c) \to \QCoh(\mathcal{O}_U),\quad
(\mathcal{F}, \beta) \mapsto \mathcal{F}
$$
と記し、補題で構成した右随伴を
$$
G : \QCoh(\mathcal{O}_U) \to \QCoh(U, R, s, t, c),\quad
\mathcal{G} \mapsto (s_*t^*\mathcal{G}, \alpha)
$$
と記す。このとき、随伴の単位 $\eta:\text{id}\to G\circ F$ を
$(\mathcal{F},\beta)$ で評価したものは写像
$$
\mathcal{F} \to s_*s^*\mathcal{F} \xrightarrow{\beta^{-1}} s_*t^*\mathcal{F}
$$
によって与えられる。確認は省略する。
\end{remark}




\begin{lemma}
\label{spaces-groupoids-lemma-push-pull}
$S$ をスキームとする。$f:Y\to X$ を $S$ 上の代数空間の射とする。
$\mathcal{F}$ を準連接 $\mathcal{O}_X$-加群、$\mathcal{G}$ を準連接
$\mathcal{O}_Y$-加群とし、$\varphi:\mathcal{G}\to f^*\mathcal{F}$ を加群写像とする。
次を仮定する。
\begin{enumerate}
\item $\varphi$ は単射である。
\item $f$ は準コンパクト、準分離的、平坦、かつ全射である。
\item $X$、$Y$ は局所 Noether 的である。
\item $\mathcal{G}$ は連接 $\mathcal{O}_Y$-加群である。
\end{enumerate}
このとき、引戻しとして定義される $\mathcal{F}\cap f_*\mathcal{G}$
$$
\xymatrix{
\mathcal{F} \ar[r] & f_*f^*\mathcal{F} \\
\mathcal{F} \cap f_*\mathcal{G} \ar[u] \ar[r] &
f_*\mathcal{G} \ar[u]
}
$$
は連接 $\mathcal{O}_X$-加群である。
\end{lemma}

\begin{proof}
Cohomology of Spaces, Lemma \ref{spaces-cohomology-lemma-coherent-Noetherian}
による連接加群の特徴づけと、連接加群が $\QCoh(\mathcal{O}_X)$ の Serre 部分圏をなすことを
自由に用いる。後者については Cohomology of Spaces, Lemma
\ref{spaces-cohomology-lemma-coherent-Noetherian-quasi-coherent-sub-quotient} を参照せよ。
$f$ が切断 $\sigma$ をもつならば、$\mathcal{F}\cap f_*\mathcal{G}$ は
$\sigma^*\mathcal{G}\to\sigma^*f^*\mathcal{F}=\mathcal{F}$ の像に含まれるので連接である。
一般の場合、$\mathcal{F}\cap f_*\mathcal{G}$ が連接であることを示すには、
% P09-ERR-1303
$f^*(\mathcal{F}\cap f_*\mathcal{G})$ が連接であることを示せば十分である
（Descent on Spaces, Lemma \ref{spaces-descent-lemma-finite-type-descends} を参照せよ）。
$f$ は平坦なので、これは $f^*\mathcal{F}\cap f^*f_*\mathcal{G}$ に等しい。
$f$ は平坦、準コンパクト、かつ準分離的なので、
$f^*f_*\mathcal{G}=p_*q^*\mathcal{G}$ である。ここで
$p,q:Y\times_X Y\to Y$ は射影である。Cohomology of Spaces, Lemma
\ref{spaces-cohomology-lemma-flat-base-change-cohomology} を参照せよ。
$p$ は切断をもつので結論が従う。
\end{proof}

\noindent
Section \ref{spaces-groupoids-section-notation} のように $B\to S$ とする。
$(U,R,s,t,c)$ を $B$ 上の代数空間における亜群とする。
$U$ は局所 Noether 的であると仮定する。以下の補題では、
準連接層 $(\mathcal{F},\alpha)$ で $(U,R,s,t,c)$ 上にあるものが {\it 連接}であるとは、
$\mathcal{F}$ が連接 $\mathcal{O}_U$-加群であることをいう。

\begin{lemma}
\label{spaces-groupoids-lemma-colimit-coherent}
Section \ref{spaces-groupoids-section-notation} のように $B\to S$ とする。
$(U,R,s,t,c)$ を $B$ 上の代数空間における亜群とする。次を仮定する。
\begin{enumerate}
\item $U$、$R$ は Noether 的である。
\item $s,t$ は平坦、準コンパクト、かつ準分離的である。
\end{enumerate}
このとき、任意の準連接加群 $(\mathcal{F},\alpha)$ で $(U,R,s,t,c)$ 上にあるものは、
連接加群のフィルター付き余極限である。
\end{lemma}

\begin{proof}
局所 Noether 的代数空間上の連接加群については、Cohomology of Spaces, Lemma
\ref{spaces-cohomology-lemma-coherent-Noetherian} の特徴づけを以後断りなく用いる。
$\mathcal{F}=\colim\mathcal{H}_i$ と、連接部分加群
$\mathcal{H}_i\subset\mathcal{F}$ のフィルター付き余極限として書ける。
Cohomology of Spaces, Lemma
\ref{spaces-cohomology-lemma-directed-colimit-coherent} を参照せよ。
$\mathcal{H}$ を $U$ 上の準連接層とする。Lemma
\ref{spaces-groupoids-lemma-construct-quasi-coherent} の準連接層 $(s_*t^*\mathcal{H},\alpha)$ で
$(U,R,s,t,c)$ 上にあるものを考える。随伴写像
$(\mathcal{F},\beta)\to(s_*t^*\mathcal{F},\alpha)$ を $\QCoh(U,R,s,t,c)$ において考える。
Remark \ref{spaces-groupoids-remark-adjunction-map} を参照せよ。次のようにおき、
$$
(\mathcal{F}_i, \beta_i) =
(\mathcal{F}, \beta)
\times_{(s_*t^*\mathcal{F}, \alpha)}
(s_*t^*\mathcal{H}_i, \alpha)
$$
$\QCoh(U,R,s,t,c)$ におけるものとする。$U$ への制限は、Lemma \ref{spaces-groupoids-lemma-abelian} の証明により
$\QCoh(U,R,s,t,c)$ 上の完全関手なので、引戻し図式
$$
\xymatrix{
\mathcal{F} \ar[r] & s_*t^*\mathcal{F} \\
\mathcal{F}_i \ar[r] \ar[u] & s_*t^*\mathcal{H}_i \ar[u]
}
$$
を得る。言い換えると $\mathcal{F}_i=\mathcal{F}\cap s_*t^*\mathcal{H}_i$ である。
Remark \ref{spaces-groupoids-remark-adjunction-map} の随伴写像の記述により、この図式は図式
$$
\xymatrix{
\mathcal{F} \ar[r] & s_*s^*\mathcal{F} \\
\mathcal{F}_i \ar[r] \ar[u] & s_*t^*\mathcal{H}_i \ar[u]
}
$$
と同型である。ここで右の鉛直矢印は、$s_*$ を写像
$$
t^*\mathcal{H}_i \to t^*\mathcal{F} \xrightarrow{\beta} s^*\mathcal{F}
$$
に適用した結果である。この矢印は $t$ が平坦なので単射である。
Lemma \ref{spaces-groupoids-lemma-push-pull} により $\mathcal{F}_i$ は連接である。
最後に、$s$ は準コンパクトかつ準分離的なので $s_*$ は余極限と可換である
（Cohomology of Schemes, Lemma \ref{coherent-lemma-colimit-cohomology} を参照せよ）。
したがって $s_*t^*\mathcal{F}=\colim s_*t^*\mathcal{H}_i$ であり、ゆえに望みどおり
$(\mathcal{F},\beta)=\colim(\mathcal{F}_i,\beta_i)$ である。
\end{proof}









\section{準連接層の結晶}
\label{spaces-groupoids-section-crystals}

\noindent
% P09-ERR-1304
$(I,\Phi,j)$ を、集合 $I$ と前関係 $j:\Phi\to I\times I$ からなる組とする。
すべての $i\in I$ に対してスキーム $X_i$ が、すべての $\phi\in\Phi$ に対して
スキームの射 $f_\phi:X_{i'}\to X_i$ で $j(\phi)=(i,i')$ となるものが与えられていると仮定する。
$$
X=(\{X_i\}_{i\in I},\{f_\phi\}_{\phi\in\Phi})
$$
とおく。{\it $X$ 上の準連接加群の結晶}とは、すべての
% P09-ERR-1305
$i\in\Ob(\mathcal{I})$ に準連接層 $\mathcal{F}_i$ で $X_i$ 上にあるものを対応させ、すべての
$\phi\in\Phi$ で $j(\phi)=(i,i')$ となるものに同型
$$
\alpha_\phi:f_\phi^*\mathcal{F}_i\longrightarrow\mathcal{F}_{i'}
$$
で $X_{i'}$ 上の準連接層のものを対応させる規則をいう。これらの準連接加群の結晶は加法圏
$\textit{CQC}(X)$ をなす\footnote{三つ組 $\phi,\phi',\phi''\in\Phi$ で
$j(\phi)=(i,i')$、$j(\phi')=(i',i'')$、$j(\phi'')=(i,i'')$ であり、
$f_{\phi''}=f_\phi\circ f_{\phi'}$ となるものの集合を特に選び、これらの三つ組について
$\alpha_{\phi'}\circ f_{\phi'}^*\alpha_\phi=\alpha_{\phi''}$ を要求することもできる。
これは加法部分圏を定める。例えば、データ $(I,\Phi)$ は添字圏の対象と矢印の集合であり得て、
$X$ はこの添字圏上のスキームの図式であり得る。Lemma
\ref{spaces-groupoids-lemma-crystals-in-quasi-coherent-modules} の結果から、この部分圏における対応する結果が
直ちに従う。}。この圏は余極限をもつ（証明は Lemma \ref{spaces-groupoids-lemma-colimits} の証明と同じである）。
すべての射 $f_\phi$ が平坦ならば、$\textit{CQC}(X)$ は Abel 圏である
（証明は Lemma \ref{spaces-groupoids-lemma-abelian} の証明と同じである）。
$\kappa$ を基数とする。準連接加群の結晶 $\mathcal{F}$ で $X$ 上にあるものが
{\it $\kappa$-生成}であるとは、各 $\mathcal{F}_i$ が $\kappa$-生成であることをいう
（Properties, Definition \ref{properties-definition-kappa-generated} を参照せよ）。

\begin{lemma}
\label{spaces-groupoids-lemma-crystals-in-quasi-coherent-modules}
上の状況で、すべての射 $f_\phi$ が平坦ならば、基数 $\kappa$ であって、
すべての対象 $(\{\mathcal{F}_i\}_{i\in I},\{\alpha_\phi\}_{\phi\in\Phi})$ で
$\textit{CQC}(X)$ のものがその
$\kappa$-生成部分加群の有向余極限となるものが存在する。
\end{lemma}

\begin{proof}
補題およびこの証明において、
$(\{\mathcal{F}_i\}_{i\in I},\{\alpha_\phi\}_{\phi\in\Phi})$ の
{\it 部分加群}とは、準連接部分加群 $\mathcal{G}_i\subset\mathcal{F}_i$ をすべての $i$ に
対して与え、等式 $\alpha_\phi(f_\phi^*\mathcal{G}_i)=\mathcal{G}_{i'}$ が
$\mathcal{F}_{i'}$ の部分層としてすべての $\phi\in\Phi$ に対して成り立つもののデータをいう。
これは、$f_\phi$ が平坦なので引戻し $f^*_\phi$ が完全、すなわち部分層を保つことから意味をもつ。
証明は Properties, Lemma \ref{properties-lemma-colimit-kappa} の証明の変形である。
まずそちらの証明を読むことを強く勧める。

\medskip\noindent
すべてのスキーム $X_i$ が affine である場合に補題を証明すれば十分であると主張する。
これを見るため、
$$
J=\coprod\nolimits_{i\in I}\{U\subset X_i\text{ affine open}\}
$$
とおき、
\begin{align*}
\Psi = & \coprod\nolimits_{\phi \in \Phi}
\{
(U, V) \mid
% P09-ERR-1308
U \subset X_i, V \subset X_{i'}\text{ affine open with } f_\phi(U) \subset V
\} \\
&
\amalg \coprod\nolimits_{i \in I}
\{
(U, U') \mid
U, U' \subset X_i\text{ affine open with } U \subset U'
\}
\end{align*}
とし、これに明らかな写像 $\Psi\to J\times J$ を備える。このとき、$(\mathcal{F},\alpha)$ は
準連接層の結晶 $(\{\mathcal{H}_j\}_{j\in J},\{\beta_\psi\}_{\psi\in\Psi})$ で
$Y=(J,\Psi)$ 上にあるものを誘導する。具体的には、
$\mathcal{H}_{(i,U)}=\mathcal{F}_i|_U$ と $(i,U)\in J$ に対しておき、
$\beta_\psi$ を $\psi\in\Psi$ に対して次のようにおく。
% P09-ERR-1309
$\alpha_\phi$ の $U$ への制限とするのは $\psi=(\phi,U,V)$ の場合であり、
$\text{id}:(\mathcal{F}_i|_{U'})|_U\to\mathcal{F}_i|_U$ とするのは
$\psi=(i,U,U')$ の場合である。
さらに、$(\{\mathcal{H}_j\}_{j\in J},\{\beta_\psi\}_{\psi\in\Psi})$ の部分加群は、
$1$ 対 $1$ に、$(\{\mathcal{F}_i\}_{i\in I},\{\alpha_\phi\}_{\phi\in\Phi})$ の
部分加群と対応する。
証明は省略する（ヒント：Sheaves, Section \ref{sheaves-section-bases} を用いよ）。
また、$\kappa$ が $Y$ に対して有効ならば、同じ $\kappa$ が $X$ に対して有効であることは
$\kappa$-生成加群の定義から明らかである。したがって、
% P09-ERR-1306
$Y$ 上の準連接層の結晶について補題を証明すれば十分である。

\medskip\noindent
すべてのスキーム $X_i$ が affine であると仮定する。$\kappa$ を、$I$ または $\Phi$ の濃度より
大きい無限基数とする。$(\{\mathcal{F}_i\}_{i\in I},\{\alpha_\phi\}_{\phi\in\Phi})$ を
$\textit{CQC}(X)$ の対象とする。各 $i$ に対して
$X_i=\Spec(A_i)$ および $M_i=\Gamma(X_i,\mathcal{F}_i)$ と書く。
すべての $\phi\in\Phi$ で $j(\phi)=(i,i')$ となるものに対して、$\alpha_\phi$ は
$A_{i'}$-加群同型
$$
\alpha_\phi:M_i\otimes_{A_i}A_{i'}\longrightarrow M_{i'}
$$
に移る。選択公理を用いて規則
% P09-ERR-1307
$$
(\phi,m)\longmapsto S(\phi,m')
$$
を選ぶ。ここで始域は、対 $(\phi,m')$ であって、$\phi\in\Phi$ かつ
$j(\phi)=(i,i')$ であり $m'\in M_{i'}$ となるものの集まりであり、値は有限部分集合
$S(\phi,m')\subset M_i$ であって、
$$
m'=\alpha_\phi\left(\sum\nolimits_{m\in S(\phi,m')}m\otimes a'_m\right)
$$
がある $a'_m\in A_{i'}$ に対して成り立つものとする。

\medskip\noindent
これらを選んだ後、任意の $\mathcal{F}_i$ の任意の $X_i$ 上の切断は、
$\kappa$-生成部分加群に含まれると主張する。これを見るため、部分集合の族
$\mathcal{S}=\{S_i\}_{i\in I}$ であって各 $S_i\subset M_i$ の濃度が高々 $\kappa$ であるものが
与えられたとする。新しい族 $\mathcal{S}'=\{S'_i\}_{i\in I}$ を
$$
S'_i = S_i \cup
\bigcup\nolimits_{(\phi, m'),\ j(\phi) = (i, i'),\ m' \in S_{i'}} S(\phi, m')
$$
によって定める。各 $S'_i$ の濃度も高々 $\kappa$ であることに注意せよ。
$\mathcal{S}^{(0)}=\mathcal{S}$、$\mathcal{S}^{(1)}=\mathcal{S}'$ とおき、帰納的に
$\mathcal{S}^{(n+1)}=(\mathcal{S}^{(n)})'$ とおく。さらに
$S_i^{(\infty)}=\bigcup_{n\geq0}S_i^{(n)}$ および
$\mathcal{S}^{(\infty)}=\{S_i^{(\infty)}\}_{i\in I}$ とおく。
構成により、すべての $\phi\in\Phi$ で $j(\phi)=(i,i')$ となるものと、
すべての $m'\in S^{(\infty)}_{i'}$ に対して、$m'$ を像
$\alpha_\phi(m\otimes1)$ の有限線形結合として書ける。ここで
$m\in S_i^{(\infty)}$ である。したがって、$N_i$ を $A_i$-部分加群であって
$M_i$ に含まれ、$S_i^{(\infty)}$ が生成するものとすると、対応する準連接部分加群
$\widetilde{N_i}\subset\mathcal{F}_i$ は $\kappa$-生成部分加群をなす。これで証明が終わる。
\end{proof}

\begin{lemma}
\label{spaces-groupoids-lemma-set-generators}
Section \ref{spaces-groupoids-section-notation} のように $B\to S$ とする。
$(U,R,s,t,c)$ を $B$ 上の代数空間における亜群とする。
$s$、$t$ が平坦ならば、集合 $T$ と対象の族
$(\mathcal{F}_t,\alpha_t)_{t\in T}$ で $\QCoh(U,R,s,t,c)$ のものが存在し、
すべての対象 $(\mathcal{F},\alpha)$ が、
対象 $(\mathcal{F}_t,\alpha_t)$ の一つと同型なその部分加群の有向余極限となるものが存在する。
\end{lemma}

\begin{proof}
この補題は、スキームにおける亜群の場合を扱う Groupoids, Lemma
\ref{groupoids-lemma-colimit-kappa} の一般化である。同じ議論をそのまま用いることはできないので、
上で展開した「準連接層の結晶」に関する内容を用いる。

\medskip\noindent
スキーム $W$ と全射 \'etale 射 $W\to U$ を選ぶ。スキーム $V$ と全射 \'etale 射
$V\to W\times_{U,s}R$ を選ぶ。スキーム $V'$ と全射 \'etale 射
$V'\to R\times_{t,U}W$ を選ぶ。スキームの集まり
$$
I = \{W, W \times_U W, V, V', V \times_R V'\}
$$
とスキームの射の集合
$$
\Phi = \{\text{pr}_i : W \times_U W \to W, V \to W, V' \to W,
V \times_R V' \to V, V \times_R V' \to V'\}
$$
を考える。$X=(I,\Phi)$ とおく。圏 $\textit{CQC}(X)$ で
$X$ 上の準連接層の結晶からなるものを定義したことを思い出そう。関手
$$
\QCoh(U, R, s, t, c) \longrightarrow \textit{CQC}(X)
$$
が存在する。これは $(\mathcal{F},\alpha)$ に、層 $\mathcal{F}|_W$ で $W$ 上にあるもの、
層 $\mathcal{F}|_{W\times_U W}$ で $W\times_U W$ 上にあるもの、
$\mathcal{F}$ の $V\to W\times_{U,s}R\to W\to U$ による引戻しで $V$ 上にあるもの、
$\mathcal{F}$ の $V'\to R\times_{t,U}W\to W\to U$ による引戻しで $V'$ 上にあるもの、
最後に $\mathcal{F}$ の
$V\times_R V'\to V\to W\times_{U,s}R\to W\to U$ による引戻しで
$V\times_R V'$ 上にあるものを対応させる。
比較写像 $\{\alpha_\phi\}_{\phi\in\Phi}$ には、引戻しの結合律から来る明らかな写像を用いるが、
写像 $\phi=\text{pr}_{V'}:V\times_R V'\to V'$ に対しては、
$\alpha:t^*\mathcal{F}\to s^*\mathcal{F}$ の $V\times_R V'$ への引戻しを用いる。
これは次の可換図式により意味をもつ。
$$
\xymatrix{
& V \times_R V' \ar[ld] \ar[rd] \\
V \ar[rd] \ar[dd] & & V' \ar[ld] \ar[dd] \\
& R \ar@<-1ex>[dd]_s \ar@<1ex>[dd]^t \\
W \ar[rd] & & W \ar[ld] \\
& U
}
$$
上で表示した関手は圏の同値ではない。しかし、$W\to U$ は全射 \'etale なので忠実である
\footnote{実際、この関手は充満忠実だが、これは必要としない。}。
上の図式のすべての射は平坦なので、これは Abel 圏の完全関手である。さらに、
$(\mathcal{F},\alpha)$ が像
$(\{\mathcal{F}_i\}_{i\in I},\{\alpha_\phi\}_{\phi\in\Phi})$ をもつとき、
$1$ 対 $1$ の対応が $(\mathcal{F},\alpha)$ の準連接部分加群と
% P09-ERR-1310
$(\{\mathcal{F}_i\}_{i\in I},\{\alpha_\phi\}_{\phi\in\Phi})$ の間にあると主張する。
実際、$(\{\mathcal{F}_i\}_{i\in I},\{\alpha_\phi\}_{\phi\in\Phi})$ の部分加群が与えられると、
$W$ 上の部分加群と射影写像 $W\times_U W\to W$ との両立性により、その部分加群は
$\mathcal{F}$ の準連接部分加群から来ることが保証される
（Properties of Spaces, Proposition
\ref{spaces-properties-proposition-quasi-coherent} による）。さらに、
$\alpha_{\text{pr}_{V'}}$ との両立性により、この部分層は $\alpha$ と両立することが保証される
（詳細は省略する）。

\medskip\noindent
基数 $\kappa$ を Lemma \ref{spaces-groupoids-lemma-crystals-in-quasi-coherent-modules} のように
系 $X=(I,\Phi)$ に対して選ぶ。Properties, Lemma
\ref{properties-lemma-set-of-iso-classes} から、$\kappa$-生成準連接層の結晶で
$X$ 上にあるものの同型類は集合をなすことが明らかである。
したがって結論は明らかである。
\end{proof}









\section{亜群と群空間}
\label{spaces-groupoids-section-groupoids-group-spaces}

\noindent
Groupoids, Section \ref{groupoids-section-groupoids-group-schemes} と比較せよ。

\begin{lemma}
\label{spaces-groupoids-lemma-groupoid-from-action}
Section \ref{spaces-groupoids-section-notation} のように $B\to S$ とする。
$(G,m)$ を $B$ 上の群代数空間とし、単位元を $e_G$、逆元を $i_G$ とする。
$X$ を $B$ 上の代数空間とし、$a:G\times_B X\to X$ を $G$ の $X$ への
$B$ 上の作用とする。このとき、代数空間における亜群
$(U,R,s,t,c,e,i)$ で $B$ 上にあるものを次のように得る。
\begin{enumerate}
\item $U=X$、$R=G\times_B X$ とおく。
\item $s:R\to U$ を $(g,x)\mapsto x$ とおく。
\item $t:R\to U$ を $(g,x)\mapsto a(g,x)$ とおく。
\item $c:R\times_{s,U,t}R\to R$ を
$((g,x),(g',x'))\mapsto(m(g,g'),x')$ とおく。
\item $e:U\to R$ を $x\mapsto(e_G(x),x)$ とおく。
\item $i:R\to R$ を $(g,x)\mapsto(i_G(g),a(g,x))$ とおく。
\end{enumerate}
\end{lemma}

\begin{proof}
省略する。ヒント：集合の水準でこれが成り立つことを示せば十分である。そのため、上の記述で
$g$ を $v$ から $a(g,v)$ への矢印とみなしたものを用いよ。
\end{proof}

\begin{lemma}
\label{spaces-groupoids-lemma-action-groupoid-modules}
Section \ref{spaces-groupoids-section-notation} のように $B\to S$ とする。
$(G,m)$ を $B$ 上の群代数空間とする。$X$ を $B$ 上の代数空間とし、
$a:G\times_B X\to X$ を $G$ の $X$ への $B$ 上の作用とする。
$(U,R,s,t,c)$ を Lemma \ref{spaces-groupoids-lemma-groupoid-from-action} で構成した代数空間における亜群とする。
規則 $(\mathcal{F},\alpha)\mapsto(\mathcal{F},\alpha)$ は、
$G$-同変 $\mathcal{O}_X$-加群の圏と $(U,R,s,t,c)$ 上の準連接加群の圏との同値を定める。
\end{lemma}

\begin{proof}
Definitions \ref{spaces-groupoids-definition-equivariant-module} および
\ref{spaces-groupoids-definition-groupoid-module} を参照すると、$t=a$ および $s=\text{pr}_1$ が
射 $R=G\times_B X\to X$ として成り立つので、主張には意味がある。
Lemma \ref{spaces-groupoids-lemma-groupoid-from-action} の翻訳を用いると、二つの定義の可換性の要請は正確に一致する。
\end{proof}





\section{固定化群代数空間}
\label{spaces-groupoids-section-stabilizer}

\noindent
Groupoids, Section \ref{groupoids-section-stabilizer} と比較せよ。
代数空間における亜群が与えられると、次のように群代数空間を得る。

\begin{lemma}
\label{spaces-groupoids-lemma-groupoid-stabilizer}
Section \ref{spaces-groupoids-section-notation} のように $B\to S$ とする。
$(U,R,s,t,c)$ を $B$ 上の代数空間における亜群とする。代数空間 $G$ をカルテシアン正方形
$$
\xymatrix{
G \ar[r] \ar[d] & R \ar[d]^{j = (t, s)} \\
U \ar[r]^-{\Delta} & U \times_B U
}
$$
によって定める。この代数空間は $U$ 上の群代数空間であり、その合成則 $m$ は合成則 $c$ から誘導される。
\end{lemma}

\begin{proof}
亜群圏では任意の対象の自己写像の集合が群をなすため、これは成り立つ。
\end{proof}

\noindent
$\Delta$ はモノ射なので、$G=j^{-1}(\Delta_{U/B})$ は $R$ の部分層である。
このように考えると、構造射 $G=j^{-1}(\Delta_{U/B})\to U$ は $s$ または $t$ のいずれからも
誘導され（両者は同じである）、$m$ は $c$ から誘導される。

\begin{definition}
\label{spaces-groupoids-definition-stabilizer-groupoid}
Section \ref{spaces-groupoids-section-notation} のように $B\to S$ とする。
$(U,R,s,t,c)$ を $B$ 上の代数空間における亜群とする。群代数空間
$j^{-1}(\Delta_{U/B})\to U$ を
{\it 代数空間における亜群 $(U,R,s,t,c)$ の固定化群}という。
\end{definition}

\noindent
文献では固定化群代数空間をしばしば $S$ と記す（おそらく stabilizer が「s」で始まるためである）。
ここでは $S$ をすでに基礎スキームに用いているので、この記法は使えない。

\begin{lemma}
\label{spaces-groupoids-lemma-groupoid-action-stabilizer}
Section \ref{spaces-groupoids-section-notation} のように $B\to S$ とする。
$(U,R,s,t,c)$ を $B$ 上の代数空間における亜群とし、$G/U$ をその固定化群とする。
$R_t/U$ を、$R$ を $U$ 上の代数空間とみなしたものと記す。ここでその構造射は
$t:R\to U$ である。標準的な左作用
$$
a : G \times_U R_t \longrightarrow R_t
$$
が存在し、これは合成則 $c$ から誘導される。
\end{lemma}

\begin{proof}
$T/B$ 上の点の言葉では、$a(g,r)=c(g,r)$ と定める。
\end{proof}








\section{亜群の制限}
\label{spaces-groupoids-section-restrict-groupoid}

\noindent
記法については Groupoids, Section \ref{groupoids-section-restrict-groupoid} を参照せよ。

\begin{lemma}
\label{spaces-groupoids-lemma-restrict-groupoid}
Section \ref{spaces-groupoids-section-notation} のように $B\to S$ とする。
$(U,R,s,t,c)$ を $B$ 上の代数空間における亜群とする。
$g:U'\to U$ を代数空間の射とする。次の図式を考える。
$$
\xymatrix{
R' \ar[d] \ar[r] \ar@/_3pc/[dd]_{t'} \ar@/^1pc/[rr]^{s'}&
R \times_{s, U} U' \ar[r] \ar[d] &
U' \ar[d]^g \\
U' \times_{U, t} R \ar[d] \ar[r] &
R \ar[r]^s \ar[d]_t &
U \\
U' \ar[r]^g &
U
}
$$
ここですべての正方形はファイバー積図式である。このとき、標準的な合成則
$c':R'\times_{s',U',t'}R'\to R'$ であって、$(U',R',s',t',c')$ が
$B$ 上の代数空間における亜群となり、$U'\to U$、$R'\to R$ が
代数空間における亜群の射
$$
(U', R', s', t', c') \to (U, R, s, t, c)
$$
で $B$ 上にあるものを定めるようなものが存在する。さらに、任意のスキーム $T$ で
$B$ 上にあるものに対し、亜群の関手
$$
(U'(T), R'(T), s', t', c') \to (U(T), R(T), s, t, c)
$$
は、$(U(T),R(T),s,t,c)$ の、写像 $U'(T)\to U(T)$ を介した制限である
（Groupoids, Section \ref{groupoids-section-restrict-groupoid} を参照せよ）。
\end{lemma}

\begin{proof}
省略する。
\end{proof}

\begin{definition}
\label{spaces-groupoids-definition-restrict-groupoid}
Section \ref{spaces-groupoids-section-notation} のように $B\to S$ とする。
$(U,R,s,t,c)$ を $B$ 上の代数空間における亜群とする。
$g:U'\to U$ を $B$ 上の代数空間の射とする。Lemma \ref{spaces-groupoids-lemma-restrict-groupoid} で構成した
代数空間における亜群の射
$(U',R',s',t',c')\to(U,R,s,t,c)$ を、
{\it $(U,R,s,t,c)$ の $U'$ への制限}という。
% P09-ERR-1311
この場合、記法 $R'=R|_{U'}$ を用いることがある。
\end{definition}

\begin{lemma}
\label{spaces-groupoids-lemma-restrict-groupoid-relation}
Definitions \ref{spaces-groupoids-definition-restrict-groupoid} および \ref{spaces-groupoids-definition-restrict-relation}
で定義した亜群と（前）同値関係の制限という概念は、Lemmas
\ref{spaces-groupoids-lemma-groupoid-pre-equivalence} および \ref{spaces-groupoids-lemma-equivalence-groupoid}
の構成を介して一致する。
\end{lemma}

\begin{proof}
ここで述べているのは、Lemma \ref{spaces-groupoids-lemma-restrict-groupoid} の $R'$ が
$$
R' = (U' \times_B U')\times_{U \times_B U} R
\longrightarrow
U' \times_B U'
$$
にも等しいということである。実際、こちらの方が補題を明確に述べる方法だったかもしれない。
\end{proof}





\section{不変部分空間}
\label{spaces-groupoids-section-invariant}

\noindent
この節では、不変部分空間の概念を簡単に論じる。

\begin{definition}
\label{spaces-groupoids-definition-invariant-open}
Section \ref{spaces-groupoids-section-notation} のように $B\to S$ とする。
$(U,R,s,t,c)$ を基礎 $B$ 上の代数空間における亜群とする。
\begin{enumerate}
\item 開部分空間 $W\subset U$ が {\it $R$-不変}であるとは、
$t(s^{-1}(W))\subset W$ を満たすことをいう。
\item 局所閉部分空間 $Z\subset U$ が、$R$ の局所閉部分空間として
$t^{-1}(Z)=s^{-1}(Z)$ を満たすとき、これを {\it $R$-不変}という。
\item 代数空間のモノ射 $T\to U$ が、$R$ 上の代数空間として
$T\times_{U,t}R=R\times_{s,U}T$ を満たすとき、これを {\it $R$-不変}という。
\end{enumerate}
\end{definition}

\noindent
開部分空間 $W\subset U$ が $R$-不変であることは、
$s^{-1}(W)=t^{-1}(W)$ を要請することとも同値である。
$W\subset U$ が $R$-不変ならば、$R$ の $W$ への制限は単に
$R_W=s^{-1}(W)=t^{-1}(W)$ である。同様に、$Z\subset U$ が
$R$-不変な局所閉部分空間ならば、$R$ の $Z$ への制限は単に
$R_Z=s^{-1}(Z)=t^{-1}(Z)$ である。

\begin{lemma}
\label{spaces-groupoids-lemma-constructing-invariant-opens}
Section \ref{spaces-groupoids-section-notation} のように $B\to S$ とする。
$(U,R,s,t,c)$ を $B$ 上の代数空間における亜群とする。
\begin{enumerate}
\item $s$ と $t$ が開ならば、任意の開部分空間 $W\subset U$ に対し、
開部分空間 $s(t^{-1}(W))$ は $R$-不変である。
\item $s$ と $t$ が開かつ準コンパクトならば、$U$ は
$R$-不変な準コンパクト開部分空間からなる開被覆をもつ。
\end{enumerate}
\end{lemma}

\begin{proof}
$s$ と $t$ が開であり、$W\subset U$ が開であると仮定する。
$s$ は開なので、$W'=s(t^{-1}(W))$ は $U$ の開部分空間である。
% P09-ERR-1312
関手的観点を用いれば、これが $R$-不変な $U$ の開部分集合であることは容易に分かるが、
有益だと考えるので、いくつかの図式を用いて直接論じる。
$t^{-1}(W')$ は射
$$
A := t^{-1}(W) \times_{s|_{t^{-1}(W)}, U, t} R
\xrightarrow{\text{pr}_1} R
$$
の像であり、$s^{-1}(W')$ は射
$$
B := R \times_{s, U, s|_{t^{-1}(W)}} t^{-1}(W)
\xrightarrow{\text{pr}_0} R.
$$
の像である。上の矢印の左側にある代数空間 $A$、$B$ は、それぞれ
$R\times_{s,U,t}R$ および $R\times_{s,U,s}R$ の開部分空間である。
Lemma \ref{spaces-groupoids-lemma-diagram} により、図式
$$
\xymatrix{
R \times_{s, U, t} R \ar[rd]_{\text{pr}_1} \ar[rr]_{(\text{pr}_1, c)} & &
R \times_{s, U, s} R \ar[ld]^{\text{pr}_0} \\
& R &
}
$$
は可換であり、横向きの矢印は同型である。さらに、
$(\text{pr}_1,c)(A)=B$ であることは明らかである。したがって
$s^{-1}(W')=t^{-1}(W')$ であり、$W'$ は $R$-不変である。これで (1) が示された。

\medskip\noindent
次に、$s$、$t$ がともに開かつ準コンパクトであると仮定する。
このとき $W\subset U$ が準コンパクト開部分空間ならば、
$W'=s(t^{-1}(W))$ も準コンパクト開部分空間であり、上の議論により不変である。
$W$ を、$U$ 上 \'etale な affine の像全体にわたって動かせば、(2) が従う。
\end{proof}





\section{商層}
\label{spaces-groupoids-section-quotient-sheaves}

\noindent
$S$ をスキームとし、$B$ を $S$ 上の代数空間とする。
$j:R\to U\times_B U$ を $B$ 上の前関係とする。各スキーム $S'$ で $S$ 上にあるものに対し、
$\sim_{S'}$ を、$j(S'):R(S')\to U(S')\times U(S')$ の像が生成する同値関係として
とることができる。したがって前層
\begin{equation}
\label{spaces-groupoids-equation-quotient-presheaf}
\begin{matrix}
(\Sch/S)^{opp}_{fppf} &
\longrightarrow &
\textit{Sets}, \\
S' &
\longmapsto  &
U(S')/\sim_{S'}
\end{matrix}
\end{equation}
を得る。$j$ は $B$ 上の代数空間の射であり、$U\times_B U$ への射でもあるので、
前層 (\ref{spaces-groupoids-equation-quotient-presheaf}) から $B$ への前層の標準的変換が存在することに注意せよ。

\begin{definition}
\label{spaces-groupoids-definition-quotient-sheaf}
$B\to S$ および前関係 $j:R\to U\times_B U$ を上のようなものとする。
この状況で {\it 商層 $U/R$} で $j$ に付随するものとは、$(\Sch/S)_{fppf}$ 上の前層
(\ref{spaces-groupoids-equation-quotient-presheaf}) の層化である。
$j:R\to U\times_B U$ が Lemma \ref{spaces-groupoids-lemma-groupoid-from-action} のような、
群代数空間 $G$ で $B$ 上にあるものの $U$ への作用から来るとき、商層を $U/G$ と記す。
\end{definition}

\noindent
これは、図式
$$
\xymatrix{
R \ar@<1ex>[r] \ar@<-1ex>[r] &
U \ar[r] &
U/R
}
$$
が $(\Sch/S)_{fppf}$ 上の集合の層の圏における余等化子図式であることを正確に意味する。
ここでも層の標準的写像 $U/R\to B$ が存在する。これは、$j$ が $B$ 上の代数空間の射であって
$U\times_B U$ への射だからである。

\begin{remark}
\label{spaces-groupoids-remark-quotient-variant}
上の構成の一つの変種として、関手
$$
\begin{matrix}
(\textit{Spaces}/B)^{opp}_{fppf} &
\longrightarrow &
\textit{Sets}, \\
X &
\longmapsto  &
U(X)/\sim_X
\end{matrix}
$$
を層化することもできた。ここで今度は、$\sim_X\subset U(X)\times U(X)$ は
$j:R(X)\to U(X)\times U(X)$ の像が生成する同値関係である。
もちろんここで $U(X)=\Mor_B(X,U)$ および $R(X)=\Mor_B(X,R)$ である。
% P09-ERR-1313
実際、（ここに Topologies of Spaces の将来の参照を挿入）の同一視を介して、結果は同じになったはずである。
\end{remark}

\begin{definition}
\label{spaces-groupoids-definition-representable-quotient}
% P09-ERR-1314
Definition \ref{spaces-groupoids-definition-quotient-sheaf} の状況において。
前関係 $j$ が {\it 代数空間で表現可能な商}をもつとは、層 $U/R$ が代数空間であることをいう。
前関係 $j$ が {\it 表現可能な商}をもつとは、層 $U/R$ がスキームによって表現されることをいう。
代数空間における亜群 $(U,R,s,t,c)$ で $B$ 上にあるものが {\it 表現可能な商}
（resp.\ {\it 代数空間で表現可能な商}）をもつとは、商 $U/R$ で $j=(t,s)$ に対するものが
表現可能（resp.\ 代数空間）であることをいう。
\end{definition}

\noindent
商 $U/R$ が $M$ によって表現されるならば（スキームまたは $S$ 上の代数空間のいずれでもよい）、
上で見たように、標準的な構造射 $M\to B$ が備わる。

\medskip\noindent
次の補題は、$M$ が商を表現することを特徴づける。
たとえば、$U\to M$ が平坦、有限表示かつ全射で、$R\cong U\times_M U$ である場合に適用できる。

\begin{lemma}
\label{spaces-groupoids-lemma-criterion-quotient-representable}
Definition \ref{spaces-groupoids-definition-quotient-sheaf} の状況において、$M$ を $S$ 上の代数空間とし、
射 $U\to M$ が存在して次を満たすと仮定する。
\begin{enumerate}
\item 射 $U\to M$ は $s,t$ を等化する。
\item 写像 $U\to M$ は層の全射である。
\item 誘導写像 $(t,s):R\to U\times_M U$ は層の全射である。
\end{enumerate}
このとき $M$ は商層 $U/R$ を表現する。
\end{lemma}

\begin{proof}
条件 (1) は $U\to M$ が $U/R$ を経由して分解することをいう。
条件 (2) は $U/R\to M$ が層の写像として全射であることをいう。
条件 (3) は $U/R\to M$ が層の写像として単射であることをいう。したがって補題が従う。
\end{proof}

\noindent
次の補題は、$j$ に前同値関係であることを要求しない場合（たとえば単なる前関係の場合）には正しくない。

\begin{lemma}
\label{spaces-groupoids-lemma-quotient-pre-equivalence}
$S$ をスキームとする。$B$ を $S$ 上の代数空間とする。
$j:R\to U\times_B U$ を $B$ 上の前同値関係とする。
$S'$ を $S$ 上のスキームとし、$a,b\in U(S')$ とする。このとき次は同値である。
\begin{enumerate}
\item $a$ と $b$ は $(U/R)(S')$ の同じ元に写る。
\item $\{f_i:S_i\to S'\}$ という $S'$ の fppf 被覆と射 $r_i:S_i\to R$ が存在し、
$a\circ f_i=s\circ r_i$ および $b\circ f_i=t\circ r_i$ が成り立つ。
\end{enumerate}
言い換えると、この場合、層の写像
$$
R \longrightarrow U \times_{U/R} U
$$
は全射である。
\end{lemma}

\begin{proof}
省略する。ヒント：これが成り立つ理由は、この場合、前層
(\ref{spaces-groupoids-equation-quotient-presheaf}) が実際に $T\mapsto U(T)/j(R(T))$ で与えられることにある。
Definition \ref{spaces-groupoids-definition-equivalence-relation} を参照すると、
$j(R(T))\subset U(T)\times U(T)$ は同値関係だからである。
\end{proof}

\begin{lemma}
\label{spaces-groupoids-lemma-quotient-pre-equivalence-relation-restrict}
$S$ をスキームとする。$B$ を $S$ 上の代数空間とする。
$j:R\to U\times_B U$ を $B$ 上の前同値関係とし、
$g:U'\to U$ を $B$ 上の代数空間の射とする。
$j':R'\to U'\times_B U'$ を $j$ の $U'$ への制限とする。商層の写像
$$
U'/R' \longrightarrow U/R
$$
は単射である。$U'\to U$ が層の写像として全射ならば、たとえば
$\{g:U'\to U\}$ が fppf 被覆ならば（Topologies on Spaces,
Definition \ref{spaces-topologies-definition-fppf-covering} を参照せよ）、
$U'/R'\to U/R$ は層の同型である。
\end{lemma}

\begin{proof}
$\xi,\xi'\in(U'/R')(S')$ を $U/R$ の同じ切断へ写る切断とする。
このとき fppf 被覆 $\mathcal{S}=\{S_i\to S'\}$ で $S'$ のものを見つけて、
$\xi|_{S_i},\xi'|_{S_i}$ が $a_i,a_i'\in U'(S_i)$ によって与えられるようにできる。
Lemma \ref{spaces-groupoids-lemma-quotient-pre-equivalence} とサイトの公理により、
% P09-ERR-1315
$\mathcal{T}$ を細分した後、射 $r_i:S_i\to R$ が存在して
$g\circ a_i=s\circ r_i$、$g\circ a_i'=t\circ r_i$ となると仮定してよい。
構成により
$R'=R\times_{U\times_S U}(U'\times_S U')$ なので、
$(r_i,(a_i,a_i'))\in R'(S_i)$ である。これは $a_i$ と $a_i'$ が
$U'/R'$ の、$S_i$ 上の同じ切断を定めることを示す。層条件により、これは
$\xi=\xi'$ を含意する。

\medskip\noindent
$U'\to U$ が層の全射ならば、$U'/R'\to U/R$ も全射である。最後に、
$\{g:U'\to U\}$ が fppf 被覆ならば、層の写像 $U'\to U$ は全射である。
Topologies on Spaces, Lemma \ref{spaces-topologies-lemma-fppf-covering-surjective} を参照せよ。
\end{proof}

\begin{lemma}
\label{spaces-groupoids-lemma-quotient-groupoid-restrict}
$S$ をスキームとする。$B$ を $S$ 上の代数空間とする。
$(U,R,s,t,c)$ を $B$ 上の代数空間における亜群とする。
% P09-ERR-1316
$g:U'\to U$ を $B$ 上の代数空間の射とする。
$(U',R',s',t',c')$ を $(U,R,s,t,c)$ の $U'$ への制限とする。商層の写像
$$
U'/R' \longrightarrow U/R
$$
は単射である。合成
$$
\xymatrix{
U' \times_{g, U, t} R \ar[r]_-{\text{pr}_1} \ar@/^3ex/[rr]^h
& R \ar[r]_s & U
}
$$
が fppf 層の全射ならば、この写像は全単射である。これはたとえば、
$\{h:U'\times_{g,U,t}R\to U\}$ が $fppf$-被覆である場合、
$U'\to U$ が層の全射である場合、または $\{g:U'\to U\}$ が fppf 位相の被覆である場合に成り立つ。
\end{lemma}

\begin{proof}
単射性は Lemmas \ref{spaces-groupoids-lemma-groupoid-pre-equivalence} および
\ref{spaces-groupoids-lemma-quotient-pre-equivalence-relation-restrict} を組み合わせると従う。
全射性を見るため（層の全射の特徴づけについては Sites, Section
\ref{sites-section-sheaves-injective} を参照せよ）、次のように論じる。
$T$ をスキームとし、$\sigma\in U/R(T)$ とする。被覆 $\{T_i\to T\}$ が存在し、
$\sigma|_{T_i}$ はある元 $f_i\in U(T_i)$ の像である。したがって、
% P09-ERR-1317
$\sigma$ は $f\in U(T)$ の像であると仮定してよい。
$h$ が層の全射であるという仮定により、fppf 被覆 $\{\varphi_i:T_i\to T\}$ と射
$f_i:T_i\to U'\times_{g,U,t}R$ を見つけて、
$f\circ\varphi_i=h\circ f_i$ とすることができる。
$f'_i=\text{pr}_0\circ f_i:T_i\to U'$ と記す。このとき
$f'_i\in U'(T_i)$ は $g\circ f'_i\in U(T_i)$ に写り、
$g\circ f'_i\sim_{T_i}h\circ f_i=f\circ\varphi_i$ であることが分かる。
記法は (\ref{spaces-groupoids-equation-quotient-presheaf}) のとおりである。すなわち、この関係を与える
$R(T_i)$ の元は $\text{pr}_1\circ f_i$ である。これは $\sigma$ の $T_i$ への制限が、
所望のとおり $U'/R'(T_i)\to U/R(T_i)$ の像に入ることを意味する。

\medskip\noindent
$\{h\}$ が fppf 被覆ならば、それは層の全射を誘導する。Topologies on Spaces,
Lemma \ref{spaces-topologies-lemma-fppf-covering-surjective} を参照せよ。
$U'\to U$ が全射ならば、$h$ も全射である。実際、$s$ は切断
% P09-ERR-1318
（すなわち亜群スキームの中立元 $e$）をもつ。
\end{proof}





\section{商スタック}
\label{spaces-groupoids-section-stacks}

\noindent
この節と続くいくつかの節では、上の Section \ref{spaces-groupoids-section-quotient-sheaves} および
Groupoids, Section \ref{groupoids-section-quotient-sheaves} の一種の一般化を記述する。
相違点は次のとおりである。商層ではなく商スタックをとる。

\medskip\noindent
$S$ をスキーム、$B$ を $S$ 上の代数空間、$(U,R,s,t,c)$ を
$B$ 上の代数空間における亜群とする。これらのデータが与えられたとき、関手
\begin{equation}
\label{spaces-groupoids-equation-quotient-stack}
\begin{matrix}
(\Sch/S)_{fppf}^{opp} &
\longrightarrow &
\textit{Groupoids} \\
S' &
\longmapsto &
(U(S'), R(S'), s, t, c)
\end{matrix}
\end{equation}
を考える。Categories, Example \ref{categories-example-functor-groupoids} により、この「亜群に値をとる前層」は
$(\Sch/S)_{fppf}$ 上の亜群にファイバー化された圏に対応する。この章ではこれを
$$
[U/_{\!p}R] \to (\Sch/S)_{fppf}
$$
と記す。下付き添字 ${}_p$ は商スタックと区別するために付けている。

\begin{definition}
\label{spaces-groupoids-definition-quotient-stack}
商スタック。$B\to S$ を上のようなものとする。
\begin{enumerate}
\item $(U,R,s,t,c)$ を $B$ 上の代数空間における亜群とする。{\it 商スタック}
$$
p : [U/R] \longrightarrow (\Sch/S)_{fppf}
$$
で $(U,R,s,t,c)$ のものとは、Stacks, Lemma
\ref{stacks-lemma-stackify-groupoids} の意味で、亜群にファイバー化された圏
$[U/_{\!p}R]$ で $(\Sch/S)_{fppf}$ 上にあり (\ref{spaces-groupoids-equation-quotient-stack}) に付随するものを
スタック化したものをいう。
\item $(G,m)$ を $B$ 上の群代数空間とする。
% P09-ERR-1319
$a:G\times_B X\to X$ を、$G$ の $B$ 上の代数空間への作用とする。{\it 商スタック}
$$
p : [X/G] \longrightarrow (\Sch/S)_{fppf}
$$
とは、Lemma \ref{spaces-groupoids-lemma-groupoid-from-action} の代数空間における亜群
$(X,G\times_B X,s,t,c)$ で $B$ 上にあるものに付随する商スタックをいう。
\end{enumerate}
\end{definition}

\noindent
したがって $[U/R]$ および $[X/G]$ は $(\Sch/S)_{fppf}$ 上の亜群に値をとるスタックである。
これらのスタックは後で非常に重要になるので、詳しく記述するのが適切である。
代数空間 $X$ で $S$ 上にあるものが与えられたとき、集合に値をとるスタックを
$\mathcal{S}_X\to(\Sch/S)_{fppf}$ と記し、これが層 $X$ に付随するものだったことを思い出そう。Categories, Lemma
\ref{categories-lemma-2-category-fibred-sets} および Stacks, Lemma
\ref{stacks-lemma-when-stack-in-sets} を参照せよ。

\begin{lemma}
\label{spaces-groupoids-lemma-quotient-stack-arrows}
$B\to S$ および $(U,R,s,t,c)$ を Definition \ref{spaces-groupoids-definition-quotient-stack} (1) のようなものとする。
標準的な $1$-射 $\pi:\mathcal{S}_U\to[U/R]$ および $[U/R]\to\mathcal{S}_B$ で、
$(\Sch/S)_{fppf}$ 上の亜群に値をとるスタックの射であるものが存在する。
合成 $\mathcal{S}_U\to\mathcal{S}_B$ は $1$-射であり、構造射 $U\to B$ に付随する。
\end{lemma}

\begin{proof}
この証明の間、亜群に値をとる前層 (\ref{spaces-groupoids-equation-quotient-stack}) に付随する
亜群にファイバー化された圏を $[U/_{\!p}R]$ と記す。スタック化の構成により、
$1$-射 $[U/_{\!p}R]\to[U/R]$ が存在する。$1$-射
$\mathcal{S}_U\to[U/R]$ は単に合成
$\mathcal{S}_U\to[U/_{\!p}R]\to[U/R]$ である。ここで最初の矢印は、スキーム $S'/S$ と
射 $x:S'\to U$ で $S$ 上にあるものに、対象 $x\in U(S')$ で
$[U/_{\!p}R]$ の $S'$ 上のファイバー圏に属するものを対応させる。

\medskip\noindent
$1$-射 $[U/R]\to\mathcal{S}_B$ を構成するには、$1$-射
$[U/_{\!p}R]\to\mathcal{S}_B$ を構成すれば十分である。Stacks, Lemma
\ref{stacks-lemma-stackify-groupoids-universal-property} を参照せよ。
$S'/S$ 上の対象には、写像
$$
U(S') \longrightarrow B(S')
$$
を用いる。これは構造射 $U\to B$ から来る。さらに明らかに、$a\in R(S')$ が始点 $s(a)\in U(S')$ と
終点 $t(a)\in U(S')$ をもつ「矢印」ならば、$s$ と $t$ は $B$ {\it 上の}射なので、
両者は同じ元 $\overline{a}$ で $B(S')$ に属するものに写る。したがって矢印 $a\in R(S')$ を
$\overline{a}$ の恒等射へ写せる。（ファイバー圏 $(\mathcal{S}_B)_{S'}$ は恒等射しか含まないので、
これは具合がよい。）この規則が、これらの分裂したファイバー圏における引戻しと両立し、
したがって所望の $1$-射 $[U/_{\!p}R]\to\mathcal{S}_B$ を定めることの確認は省略する。

\medskip\noindent
最後の主張の確認も省略する。
\end{proof}

\begin{lemma}
\label{spaces-groupoids-lemma-quotient-stack-2-arrow}
仮定と記法を Lemma \ref{spaces-groupoids-lemma-quotient-stack-arrows} のようなものとする。標準的な
$2$-射 $\alpha:\pi\circ s\to\pi\circ t$ が存在し、図式
$$
\xymatrix{
\mathcal{S}_R \ar[r]_s \ar[d]_t & \mathcal{S}_U \ar[d]^\pi \\
\mathcal{S}_U \ar[r]^-\pi & [U/R]
}
$$
を $2$-可換にする。
\end{lemma}

\begin{proof}
$S'$ を $S$ 上のスキームとする。$r:S'\to R$ を $S$ 上の射とする。
このとき $r\in R(S')$ は対象 $s\circ r,t\circ r\in U(S')$ の間の同型である。
さらに、この構成は引戻しと両立する。これにより標準的な $2$-射
$\alpha_p:\pi_p\circ s\to\pi_p\circ t$ が得られる。ここで
$\pi_p:\mathcal{S}_U\to[U/_{\!p}R]$ は Lemma \ref{spaces-groupoids-lemma-quotient-stack-arrows} の証明におけるものとする。
したがって、図式
$$
\xymatrix{
\mathcal{S}_R \ar[r]_s \ar[d]^t & \mathcal{S}_U \ar[d]^{\pi_p} \\
\mathcal{S}_U \ar[r]^-{\pi_p} & [U/_{\!p}R]
}
$$
さえ $2$-可換である。ゆえに補題の図式はなおさら $2$-可換である。
\end{proof}

\begin{remark}
\label{spaces-groupoids-remark-fundamental-square}
後の章では、集合に値をとるスタック $\mathcal{S}_X$ で $X$ に付随するものを表すのに、
単に $X$ と書く曖昧な記法を用いる。この記法を用いると、Lemma
\ref{spaces-groupoids-lemma-quotient-stack-2-arrow} の図式はよく知られた図式
$$
\xymatrix{
R \ar[r]_s \ar[d]_t & U \ar[d]^\pi \\
U \ar[r]^-\pi & [U/R]
}
$$
となる。続く節では、この図式が多くのよい性質をもつことを示す。特に、これが $2$-ファイバー積
（Section \ref{spaces-groupoids-section-quotient-stack-2-cartesian}）であり、また $2$-余等化子
（対象となる射は $s$ と $t$）に近いこと（Section \ref{spaces-groupoids-section-quotient-stacks-2-coequalize}）を示す。
\end{remark}




\section{商スタックの関手性}
\label{spaces-groupoids-section-functoriality-quotient-stacks}

\noindent
代数空間における亜群の射は、商スタックの付随する射を与える。

\begin{lemma}
\label{spaces-groupoids-lemma-quotient-stack-functorial}
$S$ をスキームとする。$B$ を $S$ 上の代数空間とする。
$f:(U,R,s,t,c)\to(U',R',s',t',c')$ を $B$ 上の代数空間における亜群の射とする。
このとき $f$ は商スタックの標準的な $1$-射
$$
[f] : [U/R] \longrightarrow [U'/R'].
$$
を誘導する。
\end{lemma}

\begin{proof}
$[U/_{\!p}R]$ および $[U'/_{\!p}R']$ を、関手 (\ref{spaces-groupoids-equation-quotient-stack}) に付随する
基礎サイト $(\Sch/S)_{fppf}$ 上の亜群にファイバー化された圏と記す。
$f$ が $1$-射 $[U/_{\!p}R]\to[U'/_{\!p}R']$ を定めることは明らかであり、これを
$[U'/R']$ のスタック化関手と合成して $[U/_{\!p}R]\to[U'/R']$ を得る。
% P09-ERR-1320
次に、スタック化関手 $[U/_{\!p}R]\to[U/R]$ の普遍性、Stacks, Lemma
\ref{stacks-lemma-stackify-groupoids-universal-property} により、$[U/R]\to[U'/R']$ を得る。
\end{proof}

\noindent
$B\to S$ および $f:(U,R,s,t,c)\to(U',R',s',t',c')$ を Lemma
\ref{spaces-groupoids-lemma-quotient-stack-functorial} のようなものとする。この状況で、$B$ 上の
代数空間における第3の亜群を次のように定める。ここでは $T$-値点の言葉を用い、
$T$ は $B$ 上の（動く）スキームとする。
\begin{enumerate}
\item $U''=U\times_{f,U',t'}R'$ とおく。したがって $T$-値点は
対 $(u,r')$ で $f(u)=t'(r')$ を満たすものである。
\item $R''=R\times_{f\circ s,U',t'}R'$ とおく。したがって $T$-値点は
対 $(r,r')$ で $f(s(r))=t'(r')$ を満たすものである。
\item $s'':R''\to U''$ を $s''(r,r')=(s(r),r')$ によって与える。
\item $t'':R''\to U''$ を $t''(r,r')=(t(r),c'(f(r),r'))$ によって与える。
\item $c'':R''\times_{s'',U'',t''}R''\to R''$ を
$c''((r_1,r'_1),(r_2,r'_2))=(c(r_1,r_2),r'_2)$ によって与える。
\end{enumerate}
$c''$ の式には意味がある。実際、$s''(r_1,r'_1)=t''(r_2,r'_2)$ である。
$c''$ が結合的であることは明らかである。恒等元 $e''$ は
$e''(u,r)=(e(u),r)$ によって与えられる。$(r,r')$ の逆元は
$(i(r),c'(f(r),r'))$ によって与えられる。したがって実際に $B$ 上の
代数空間における亜群を得る。

\medskip\noindent
写像 $U''\to U$ および $R''\to R$ が、代数空間における亜群の射
$g:(U'',R'',s'',t'',c'')\to(U,R,s,t,c)$ で $B$ 上にあるものを定めることは明らかである。
さらに、写像 $U''\to U'$、$(u,r')\mapsto s'(r')$ および
$R''\to U'$、$(r,r')\mapsto s'(r')$ は、実際
$(U'',R'',s'',t'',c'')$ が $U'$ 上の代数空間における亜群であることを示す。

\begin{lemma}
\label{spaces-groupoids-lemma-cartesian-square-of-morphism}
記法と仮定を Lemma \ref{spaces-groupoids-lemma-quotient-stack-functorial} のようなものとする。
$(U'',R'',s'',t'',c'')$ を、上で構成した $B$ 上の代数空間における亜群とする。
$2$-可換正方形
$$
\xymatrix{
[U''/R''] \ar[d] \ar[r]_{[g]} & [U/R] \ar[d]^{[f]} \\
\mathcal{S}_{U'} \ar[r] & [U'/R']
}
$$
が存在し、これは $[U''/R'']$ を $2$-ファイバー積と同一視する。
\end{lemma}

\begin{proof}
写像 $[f]$ および $[g]$ は Lemma \ref{spaces-groupoids-lemma-quotient-stack-functorial} の適用から来て、
他の二つの写像は Lemma \ref{spaces-groupoids-lemma-quotient-stack-arrows} から来る
（さらに $(U'',R'',s'',t'',c'')$ は $U'$ 上にある）。
$2$-ファイバー積の性質を示すには、亜群にファイバー化された圏の図式
$$
\xymatrix{
[U''/_{\!p}R''] \ar[d] \ar[r]_{[g]} & [U/_{\!p}R] \ar[d]^{[f]} \\
\mathcal{S}_{U'} \ar[r] & [U'/_{\!p}R']
}
$$
について補題を証明すれば十分である。Stacks, Lemma
\ref{stacks-lemma-stackification-fibre-product-categories-fibred-in-groupoids} を参照せよ。
言い換えると、$2$-ファイバー積
% P09-ERR-1321
$\mathcal{S}_U\times_{[U'/_{\!p}R']}[U/_{\!p}R]$ で $T$ 上にあるものの対象が、
$T$-値点で $U''$ のものに対応し、
射についても同様であることを示せば十分である。もちろん、これはまさに最初に
$U''$ および $R''$ を構成した方法である。

\medskip\noindent
詳しく述べると、$\mathcal{S}_U\times_{[U'/_{\!p}R']}[U/_{\!p}R]$ の $T$ 上の対象は
三つ組 $(u',u,r')$ である。ここで $u'$ は $T$-値点で $U'$ のもの、$u$ は
$T$-値点で $U$ のもの、$r'$ は $u'$ から $f(u)$ への射で $[U'/R']_T$ におけるもの、すなわち
% P09-ERR-1322
$r'$ は $T$-値点で $R$ のものであり、$s'(r')=u'$ および $t'(r')=f(u)$ を満たす。
$u'$ を忘れても情報は失われず、これらの対象が
% P09-ERR-1323
$T$-値点で $R''$ のものと一対一に対応することが分かる。

\medskip\noindent
射についても同様である。$(u'_1,u_1,r'_1)$ および $(u'_2,u_2,r'_2)$ を
$T$ 上のファイバー積の二つの対象とする。このとき $(u'_2,u_2,r'_2)$ から
$(u'_1,u_1,r'_1)$ への射は $(1,r)$ によって与えられる。ここで
$1:u'_1\to u'_2$ は単に $u'_1=u'_2$ を意味し
% P09-ERR-1324
（これは $\mathcal{S}_U$ が集合にファイバー化されているためである）、$r$ は
$T$-値点で $R$ のものであり、$s(r)=u_2$、$t(r)=u_1$、さらに
$c'(f(r),r'_2)=r'_1$ を満たすものである。したがって矢印
$$
(1, r) : (u'_2, u_2, r'_2) \to (u'_1, u_1, r'_1)
$$
は対 $(r,r'_2)$ を知れば完全に決まる。ゆえに矢印の関手は $R''$ によって表現され、さらに
射 $s''$、$t''$、$c''$ は明らかに $2$-ファイバー積
$\mathcal{S}_U\times_{[U'/_{\!p}R']}[U/_{\!p}R]$ における始点、終点、合成に対応する。
\end{proof}







\section{商スタックの 2-カルテシアン正方形}
\label{spaces-groupoids-section-quotient-stack-2-cartesian}

\noindent
この節では、商スタックの $\mathit{Isom}$-層を計算し、商スタックを定義する図式が
$2$-ファイバー積であることを導く。

\begin{lemma}
\label{spaces-groupoids-lemma-quotient-stack-morphisms}
$B\to S$、$(U,R,s,t,c)$、および $\pi:\mathcal{S}_U\to[U/R]$ を
Lemma \ref{spaces-groupoids-lemma-quotient-stack-arrows} のようなものとする。
$S'$ を $S$ 上のスキームとする。$x,y\in\Ob([U/R]_{S'})$ を $S'$ 上の商スタックの対象とする。
$x=\pi(x')$ および $y=\pi(y')$ となる射 $x',y':S'\to U$ があるならば、
$$
\mathit{Isom}(x, y) = S' \times_{(y', x'), U \times_S U} R
$$
が $S'$ 上の層として成り立つ。
\end{lemma}

\begin{proof}
$[U/_{\!p}R]$ を、亜群に値をとる前層 (\ref{spaces-groupoids-equation-quotient-stack}) に付随する
亜群にファイバー化された圏とする。これは Lemma \ref{spaces-groupoids-lemma-quotient-stack-arrows} の証明におけるものと同じである。
構成により、層 $\mathit{Isom}(x,y)$ は前層 $\mathit{Isom}(x',y')$ に付随する層である。
一方、$[U/_{\!p}R]$ における射の定義により
$$
\mathit{Isom}(x', y') = S' \times_{(y', x'), U \times_S U} R
$$
であり、右辺は代数空間、したがって層である。
\end{proof}

\begin{lemma}
\label{spaces-groupoids-lemma-quotient-stack-2-cartesian}
$B\to S$、$(U,R,s,t,c)$、および $\pi:\mathcal{S}_U\to[U/R]$ を
Lemma \ref{spaces-groupoids-lemma-quotient-stack-arrows} のようなものとする。
Lemma \ref{spaces-groupoids-lemma-quotient-stack-2-arrow} の $2$-可換正方形
$$
\xymatrix{
\mathcal{S}_R \ar[r]_s \ar[d]_t & \mathcal{S}_U \ar[d]^\pi \\
\mathcal{S}_U \ar[r]^-\pi & [U/R]
}
$$
は $2$-ファイバー積であり、$(\Sch/S)_{fppf}$ の亜群に値をとるスタックのものとして考える。
\end{lemma}

\begin{proof}
Stacks, Lemma \ref{stacks-lemma-2-product-stacks-in-groupoids} により、この補題には意味がある。
同じ補題は、関手
$$
\mathcal{S}_R \longrightarrow \mathcal{S}_U \times_{[U/R]} \mathcal{S}_U
$$
で $r:T\to R$ を $(T,t(r),s(r),\alpha(r))$ に写すものが同値であることを示す必要があるとも述べている。ここで右辺は Categories, Lemma
\ref{categories-lemma-2-product-categories-over-C} に記述された $2$-ファイバー積である。
定義を書き下せば、これはちょうど Lemma \ref{spaces-groupoids-lemma-quotient-stack-morphisms} の内容である。
% P09-ERR-1325
（別証明：この状況で Lemma \ref{spaces-groupoids-lemma-cartesian-square-of-morphism} の意味を書き下すことも、
結果を与える。）
\end{proof}

\begin{lemma}
\label{spaces-groupoids-lemma-quotient-stack-isom}
$B\to S$ および $(U,R,s,t,c)$ を Definition \ref{spaces-groupoids-definition-quotient-stack} (1) のようなものとする。
$T$ を $S$ 上の任意のスキームとし、$x,y$ を $[U/R]$ の $T$ 上の対象とする。
層 $\mathit{Isom}(x,y)$ で $(\Sch/T)_{fppf}$ 上にあるものは次の性質をもつ。fppf 被覆
$\{T_i\to T\}_{i\in I}$ が存在し、
$\mathit{Isom}(x,y)|_{(\Sch/T_i)_{fppf}}$ は代数空間によって表現可能である。
\end{lemma}

\begin{proof}
Lemma \ref{spaces-groupoids-lemma-quotient-stack-morphisms} と、商スタックの構成により $x$ と $y$ の双方が
fppf 位相で局所的に $\mathcal{S}_U$ の対象から来るという事実から直ちに従う。
\end{proof}














\section{商スタックの 2-余等化子性質}
\label{spaces-groupoids-section-quotient-stacks-2-coequalize}

\noindent
亜群には合成があり、これは上の補題の標準的 $2$-射に対するコサイクル条件を導く。
正確な定式化を与えるため、Categories, Sections
\ref{categories-section-formal-cat-cat} および \ref{categories-section-2-categories}
で導入された記法を用いる。

\begin{lemma}
\label{spaces-groupoids-lemma-quotient-stack-cocycle}
仮定と記法を Lemmas \ref{spaces-groupoids-lemma-quotient-stack-arrows} および
\ref{spaces-groupoids-lemma-quotient-stack-2-arrow} のようなものとする。次の図式の鉛直合成
$$
\xymatrix@C=15pc{
\mathcal{S}_{R \times_{s, U, t} R}
\ruppertwocell^{\pi \circ s \circ \text{pr}_1 = \pi \circ s \circ c}{\ \ \ \ \ \ \alpha \star \text{id}_{\text{pr}_1}}
\ar[r]_(.3){\pi \circ t \circ \text{pr}_1 = \pi \circ s \circ \text{pr}_0}
\rlowertwocell_{\pi \circ t \circ \text{pr}_0 = \pi \circ t \circ c}{\ \ \ \ \ \ \alpha \star \text{id}_{\text{pr}_0}}
&
[U/R]
}
$$
は $2$-射 $\alpha\star\text{id}_c$ である。式で書けば
$\alpha \star \text{id}_c =
(\alpha \star \text{id}_{\text{pr}_0})
\circ
(\alpha \star \text{id}_{\text{pr}_1})
$ である。
\end{lemma}

\begin{proof}
二つ注意する。
\begin{enumerate}
\item 式
$\alpha\star\text{id}_c=(\alpha\star\text{id}_{\text{pr}_0})\circ
(\alpha\star\text{id}_{\text{pr}_1})$ に意味があるためには、{\it 等式}
$\pi\circ s\circ\text{pr}_1=\pi\circ s\circ c$、
$\pi\circ t\circ\text{pr}_1=\pi\circ s\circ\text{pr}_0$、および
$\pi\circ t\circ\text{pr}_0=\pi\circ t\circ c$ を認識する必要がある。
すなわち、第2の等式は鉛直合成 $\circ$ に意味があることを含意し、他の二つは式の両辺が
同じ始域と終域をもつ $2$-射であることを保証する。
\item この補題が成り立つ理由は、亜群に値をとる前層
(\ref{spaces-groupoids-equation-quotient-stack}) に付随する亜群にファイバー化された圏
$[U/_{\!p}R]$ における合成が、合成則
$c:R\times_{s,U,t}R\to R$ から来ることにある。
\end{enumerate}
補題の証明は省略する。
\end{proof}

\noindent
補題の状況では、実際に等式
$s\circ\text{pr}_1=s\circ c$、
$t\circ\text{pr}_1=s\circ\text{pr}_0$、および
$t\circ\text{pr}_0=t\circ c$ が、$\pi$ と合成する前に成り立つことに注意せよ。
したがって、下の補題の式は上の補題の式とまったく同じ仕方で意味をもつ。

\begin{lemma}
\label{spaces-groupoids-lemma-quotient-stack-2-coequalizer}
仮定と記法を Lemmas \ref{spaces-groupoids-lemma-quotient-stack-arrows} および
\ref{spaces-groupoids-lemma-quotient-stack-2-arrow} のようなものとする。
Lemma \ref{spaces-groupoids-lemma-quotient-stack-2-arrow} の $2$-可換図式は、次の意味で $2$-余等化子である。
次が与えられているとする。
\begin{enumerate}
\item 亜群に値をとるスタック $\mathcal{X}$ で $(\Sch/S)_{fppf}$ 上にあるもの、
\item $1$-射 $f:\mathcal{S}_U\to\mathcal{X}$、
\item $2$-射 $\beta:f\circ s\to f\circ t$。
\end{enumerate}
さらに
$$
\beta \star \text{id}_c
=
(\beta \star \text{id}_{\text{pr}_0})
\circ
(\beta \star \text{id}_{\text{pr}_1})
$$
を満たすとする。このとき $1$-射 $[U/R]\to\mathcal{X}$ が存在して、図式
$$
\xymatrix{
\mathcal{S}_R \ar[r]_s \ar[d]^t & \mathcal{S}_U \ar[d] \ar[ddr]^f \\
\mathcal{S}_U \ar[r] \ar[rrd]_f & [U/R] \ar[rd] \\
& & \mathcal{X}
}
$$
を $2$-可換にする。
\end{lemma}

\begin{proof}
補題のような $\mathcal{X}$、$f$、$\beta$ が与えられたとする。Stacks, Lemma
\ref{stacks-lemma-stackify-groupoids-universal-property} により、$1$-射
$g:[U/_{\!p}R]\to\mathcal{X}$ を構成すれば十分である。まず、$1$-射
$\mathcal{S}_U\to[U/_{\!p}R]$ は対象上全単射であることに注意する。したがって対象上では、
$g(x)=f(x)$ とおける。ここで $x\in\Ob(\mathcal{S}_U)=\Ob([U/_{\!p}R])$ である。
射 $\varphi:x\to y$ で $[U/_{\!p}R]$ のものは可換図式
$$
\xymatrix{
S_2 \ar[dd]_h \ar[r]_x \ar[dr]_\varphi & U \\
& R \ar[u]_s \ar[d]^t \\
S_1 \ar[r]^y & U.
}
$$
から来る。したがって $g(\varphi)$ を合成
$$
\xymatrix{
f(x) \ar@{=}[r] \ar[rrrrrd] &
f(s \circ \varphi) \ar@{=}[r] &
(f \circ s)(\varphi) \ar[r]^\beta &
(f \circ t)(\varphi) \ar@{=}[r] &
f(t \circ \varphi) \ar@{=}[r] &
f(y \circ h) \ar[d] \\
& & & & & f(y).
}
$$
に等しいとおける。鉛直矢印は、関手 $f$ を標準的な射
$y\circ h\to y$ で $\mathcal{S}_U$ におけるもの（すなわち、$h$ を持ち上げ、
終点 $y$ をもつ強カルテシアン射）に適用した結果である。
% P09-ERR-1326
このように定めた $f$ が、少なくともファイバー圏上で合成と両立することを確認しよう。
$S'$ を $S$ 上のスキームとし、$a:S'\to R\times_{s,U,t}R$ を射とする。この状況で
$x=s\circ\text{pr}_1\circ a=s\circ c\circ a$、
$y=t\circ\text{pr}_1\circ a=s\circ\text{pr}_0\circ a$、および
% P09-ERR-1327
$z=t\circ\text{pr}_0\circ a=t\circ\text{pr}_0\circ c$ とおくと、可換図式
$$
\xymatrix{
x \ar[rr]_{c \circ a} \ar[rd]_{\text{pr}_1 \circ a} & & z \\
& y \ar[ru]_{\text{pr}_0 \circ a}
}
$$
をファイバー圏 $[U/_{\!p}R]_{S'}$ において得る。さらに、このファイバー圏の任意の可換三角形は
この形をもつ。このとき、上の定義により、$f$ がこれを可換図式に写すための必要十分条件は、図式
$$
\xymatrix{
& (f \circ s)(c \circ a) \ar[r]_-{\beta} &
(f \circ t)(c \circ a) \ar@{=}[rd] & \\
(f \circ s)(\text{pr}_1 \circ a) \ar[rd]^\beta \ar@{=}[ru] & & &
(f \circ t)(\text{pr}_0 \circ a) \\
& (f \circ t)(\text{pr}_1 \circ a) \ar@{=}[r] &
(f \circ s)(\text{pr}_0 \circ a) \ar[ru]^\beta
}
$$
が可換であることであり、これは補題の式が表す条件にほかならない。
$f$ が恒等射を恒等射に写し、任意の射の合成と両立することの確認は省略する。
\end{proof}












\section{商スタックの明示的記述}
\label{spaces-groupoids-section-explicit-quotient-stacks}

\noindent
結果を定式化するため、いくつかの記法を導入する必要がある。
$B\to S$ および $(U,R,s,t,c)$ を Definition \ref{spaces-groupoids-definition-quotient-stack} (1) のようなものとする。
$T$ を $S$ 上のスキームとする。$\mathcal{T}=\{T_i\to T\}_{i\in I}$ を fppf 被覆とする。
{\it $[U/R]$-降下データ}で $\mathcal{T}$ に関するものは系 $(u_i,r_{ij})$ によって与えられる。ここで
\begin{enumerate}
\item 各 $i$ に対し射 $u_i:T_i\to U$ があり、
\item 各 $i,j$ に対し射 $r_{ij}:T_i\times_T T_j\to R$ がある。
\end{enumerate}
さらに次を満たすものとする。
\begin{enumerate}
\item[(a)] 射 $T_i\times_T T_j\to U$ として
$$
s \circ r_{ij} = u_i \circ \text{pr}_0
\quad\text{and}\quad
t \circ r_{ij} = u_j \circ \text{pr}_1,
$$
が成り立つ。
\item[(b)] 射 $T_i\times_T T_j\times_T T_k\to R$ として
$$
c \circ (r_{jk} \circ \text{pr}_{12}, r_{ij} \circ \text{pr}_{01})
= r_{ik} \circ \text{pr}_{02}.
$$
が成り立つ。
\end{enumerate}
{\it 射 $(u_i,r_{ij})\to(u'_i,r'_{ij})$} で、二つの $[U/R]$-降下データで
同じ被覆 $\mathcal{T}$ 上にあるものの間のものとは、族 $(r_i:T_i\to R)$ で次を満たすものをいう。
\begin{enumerate}
\item[] $(\alpha)$ 射 $T_i\to U$ として
$$
u_i = s \circ r_i
\quad\text{and}\quad
u'_i = t \circ r_i
$$
が成り立つ。
\item[] $(\beta)$ 射 $T_i\times_T T_j\to R$ として
$$
c \circ (r'_{ij}, r_i \circ \text{pr}_0)
=
c \circ (r_j \circ \text{pr}_1, r_{ij}).
$$
が成り立つ。
\end{enumerate}
固定した被覆に関する降下データの射には自然な合成則があり、降下データの圏を得る。
この圏は亜群である。最後に、$\mathcal{T}'=\{T'_j\to T\}_{j\in J}$ を
$\mathcal{T}$ の細分となる第2の fppf 被覆とすると、降下データの引戻しという概念がある。
この場合、それを明示的に記述するのは特に容易である。すなわち、$\alpha:J\to I$ および
% P09-ERR-1328
$\varphi_j:T'_j\to T_{\alpha(i)}$ が被覆の射ならば、降下データ $(u_i,r_{ii'})$ の引戻しは単に
$$
(u_{\alpha(i)} \circ \varphi_j,
r_{\alpha(j)\alpha(j')} \circ \varphi_j \times \varphi_{j'}).
$$
である。このように定めた引戻しは、$\mathcal{T}$ 上の降下データの圏から
% P09-ERR-1329
$\mathcal{T}'$ 上の降下データの圏への関手を定める。

\begin{lemma}
\label{spaces-groupoids-lemma-quotient-stack-objects}
$B\to S$ および $(U,R,s,t,c)$ を Definition \ref{spaces-groupoids-definition-quotient-stack} (1) のようなものとする。
$\pi:\mathcal{S}_U\to[U/R]$ を Lemma \ref{spaces-groupoids-lemma-quotient-stack-arrows} のようなものとする。
$T$ を $S$ 上のスキームとする。
\begin{enumerate}
\item 各対象 $x$ でファイバー圏 $[U/R]_T$ のものに対し、fppf 被覆
$\{f_i:T_i\to T\}_{i\in I}$ が存在し、$f_i^*x\cong\pi(u_i)$ が
ある $u_i\in U(T_i)$ に対して成り立つ。
\item 同型の合成
$$
\pi(u_i \circ \text{pr}_0)
=
\text{pr}_0^*\pi(u_i)
\cong
\text{pr}_0^*f_i^*x
\cong
\text{pr}_1^*f_j^*x
\cong
\text{pr}_1^*\pi(u_j)
=
\pi(u_j \circ \text{pr}_1)
$$
は $\pi(r_{ij})$ の形であり、ある射 $r_{ij}:T_i\times_T T_j\to R$ に対するものである。
\item 系 $(u_i,r_{ij})$ は上で定義した $[U/R]$-降下データをなす。
\item 任意の $[U/R]$-降下データ $(u_i,r_{ij})$ はこのようにして生じる。
\item $x$ が上のように $(u_i,r_{ij})$ に対応し、
$y\in\Ob([U/R]_T)$ が $(u'_i,r'_{ij})$ に対応するならば、標準的な全単射
$$
\Mor_{[U/R]_T}(x, y)
\longleftrightarrow
\left\{
\begin{matrix}
\text{morphisms }(u_i, r_{ij}) \to (u'_i, r'_{ij})\\
\text{of }[U/R]\text{-descent data}
\end{matrix}
\right\}
$$
が存在する。
\item この対応は fppf 被覆の細分と両立する。
\end{enumerate}
\end{lemma}

\begin{proof}
% P09-ERR-1330
主張 (1) はスタック化の構成の一部である。(2) は Lemma
\ref{spaces-groupoids-lemma-quotient-stack-morphisms} から従う。(3) の確認は省略する。
(4) は、スタックではすべての降下データが有効であるという事実の翻訳である。
(5) と (6) の確認は省略する。
\end{proof}








\section{制限と商スタック}
\label{spaces-groupoids-section-quotient-stack-restrict}

\noindent
この節では、制限をとるとき商スタックに何が起こるかを調べる。

\begin{lemma}
\label{spaces-groupoids-lemma-quotient-stack-restrict}
記法と仮定を Lemma \ref{spaces-groupoids-lemma-quotient-stack-functorial} のようなものとする。商スタックの射
$$
[f] : [U/R] \longrightarrow [U'/R']
$$
が充満忠実であるための必要十分条件は、$R$ が $R'$ の制限であって射 $f:U\to U'$ を介したものであることである。
\end{lemma}

\begin{proof}
$x,y$ を $[U/R]$ の対象でスキーム $T/S$ 上にあるものとする。$x',y'$ を
$x,y$ の像で圏 $[U'/R']_T$ におけるものとする。関手 $[f]$ が充満忠実であるための必要十分条件は、層の写像
$$
\mathit{Isom}(x, y) \longrightarrow \mathit{Isom}(x', y')
$$
がすべての $T,x,y$ に対して同型であることである。これは $T$ 上局所的に（fppf 位相で）
検査できる。したがって Lemma \ref{spaces-groupoids-lemma-quotient-stack-objects} により、$x,y$ は
$a,b\in U(T)$ から来ると仮定してよい。この場合 $x',y'$ は $f\circ a,f\circ b$ に対応する。
Lemma \ref{spaces-groupoids-lemma-quotient-stack-morphisms} により、表示した層の写像はこの場合
$$
T \times_{(a, b), U \times_B U} R
\longrightarrow
T \times_{f \circ a, f \circ b, U' \times_B U'} R'.
$$
となる。$R$ が制限ならばこれは同型である。実際この場合
$R=(U\times_B U)\times_{U'\times_B U'}R'$ である。Lemma
\ref{spaces-groupoids-lemma-restrict-groupoid-relation} とその証明を参照せよ。逆に、最後に表示した写像が
すべての $T,a,b$ に対して同型ならば、
$R=(U\times_B U)\times_{U'\times_B U'}R'$ が従う。すなわち $R$ は $R'$ の制限である。
\end{proof}

\begin{lemma}
\label{spaces-groupoids-lemma-quotient-stack-restrict-equivalence}
記法と仮定を Lemma \ref{spaces-groupoids-lemma-quotient-stack-functorial} のようなものとする。商スタックの射
$$
[f] : [U/R] \longrightarrow [U'/R']
$$
が同値であるための必要十分条件は次である。
\begin{enumerate}
\item $(U,R,s,t,c)$ は $(U',R',s',t',c')$ の制限であり、$f:U\to U'$ を介したものである。
\item 写像
$$
\xymatrix{
U \times_{f, U', t'} R' \ar[r]_-{\text{pr}_1} \ar@/^3ex/[rr]^h
& R' \ar[r]_{s'} & U'
}
$$
は層の全射である。
\end{enumerate}
(2) は、たとえば $\{h:U\times_{f,U',t'}R'\to U'\}$ が fppf 被覆である場合、
$f:U\to U'$ が層の全射である場合、または $\{f:U\to U'\}$ が fppf 被覆である場合に成り立つ。
\end{lemma}

\begin{proof}
(1) が充満忠実性と同値であることは Lemma \ref{spaces-groupoids-lemma-quotient-stack-restrict} によりすでに分かっている。
% P09-ERR-1331
したがって (1) が成り立ち、$[f]$ が充満忠実であると仮定してよい。目標は、この仮定の下で
$[f]$ が同値であることと (2) が成り立つことが同値であると示すことである。同値を特徴づける
Stacks, Lemma \ref{stacks-lemma-characterize-essentially-surjective-when-ff} を用いてよい。

\medskip\noindent
(2) を仮定する。Stacks, Lemma \ref{stacks-lemma-characterize-essentially-surjective-when-ff} を用いて
$[f]$ が同値であることを証明する。$T$ をスキームとし、$x'\in\Ob([U'/R']_T)$ とする。
被覆 $\{g_i:T_i\to T\}$ が存在し、$g_i^*x'$ はある元 $a'_i\in U'(T_i)$ の像である。
Lemma \ref{spaces-groupoids-lemma-quotient-stack-objects} を参照せよ。したがって $x'$ は
$a'\in U'(T)$ の像であると仮定してよい。$h$ が層の全射であるという仮定により、
fppf 被覆 $\{\varphi_i:T_i\to T\}$ と射
% P09-ERR-1332
$b_i:T_i\to U\times_{g,U',t'}R'$ を見つけて、$a'\circ\varphi_i=h\circ b_i$ とできる。
$a_i=\text{pr}_0\circ b_i:T_i\to U$ と記す。このとき $a_i\in U(T_i)$ は
$f\circ a_i\in U'(T_i)$ に写り、$f\circ a_i\cong_{T_i}h\circ b_i=a'\circ\varphi_i$ である。
ここで $\cong_{T_i}$ はファイバー圏 $[U'/R']_{T_i}$ における同型を表す。すなわち、この同型を与える
$R'(T_i)$ の元は $\text{pr}_1\circ b_i$ である。これは
% P09-ERR-1333
$x$ の $T_i$ への制限が、所望のとおり関手 $[U/R]_{T_i}\to[U'/R']_{T_i}$ の本質像に入ることを意味する。

\medskip\noindent
$[f]$ が同値であると仮定する。$\xi'\in[U'/R']_{U'}$ を $U'$ の恒等射に対応する対象とする。
Stacks, Lemma \ref{stacks-lemma-characterize-essentially-surjective-when-ff} を適用すると、fppf 被覆
$\mathcal{U}'=\{g'_i:U'_i\to U'\}$ が存在し、$(g'_i)^*\xi'\cong[f](\xi_i)$ が
ある $\xi_i$ で $[U/R]_{U'_i}$ に属するものに対して成り立つことが分かる。被覆 $\mathcal{U}'$ を細分すれば
（Lemma \ref{spaces-groupoids-lemma-quotient-stack-objects} を用いる）、$\xi_i$ は射 $a_i:U'_i\to U$ から来ると仮定してよい。
$[f](\xi_i)\cong(g'_i)^*\xi'$ であるという事実は、必要なら被覆 $\mathcal{U}'$ をもう一度細分した後、
射 $r'_i:U'_i\to R'$ で $t'\circ r'_i=f\circ a_i$ および
$s'\circ r'_i=\text{id}_{U'}\circ g'_i$ を満たすものが存在することを意味する。図示すると
$$
\xymatrix{
U \ar[d]^f & & U'_i \ar[ll]^{a_i} \ar[ld]_{r'_i} \ar[d]^{g'_i} \\
U' & R' \ar[l]_{t'} \ar[r]^{s'} & U'
}
$$
である。したがって $(a_i,r'_i):U'_i\to U\times_{g,U',t'}R'$ は射であり、
$h\circ(a_i,r'_i)=g'_i$ を満たす。よって $\{h:U\times_{g,U',t'}R'\to U'\}$ は
fppf 被覆 $\mathcal{U}'$ によって細分できる。これは $h$ が層の全射を誘導することを意味する。
Topologies on Spaces, Lemma \ref{spaces-topologies-lemma-fppf-covering-surjective} を参照せよ。

\medskip\noindent
$\{h\}$ が fppf 被覆ならば、それは層の全射を誘導する。Topologies on Spaces, Lemma
\ref{spaces-topologies-lemma-fppf-covering-surjective} を参照せよ。
% P09-ERR-1334
$U'\to U$ が全射ならば、$h$ も全射である。実際、$s$ は切断
（すなわち代数空間における亜群の中立元 $e$）をもつ。
\end{proof}

\begin{lemma}
\label{spaces-groupoids-lemma-criterion-fibre-product}
記法と仮定を Lemma \ref{spaces-groupoids-lemma-quotient-stack-functorial} のようなものとする。図式
$$
\xymatrix{
R \ar[d]_s \ar[r]_f & R' \ar[d]^{s'} \\
U \ar[r]^f & U'
}
$$
がカルテシアンであると仮定する。このとき
$$
\xymatrix{
\mathcal{S}_U \ar[d] \ar[r] & [U/R] \ar[d]^{[f]} \\
\mathcal{S}_{U'} \ar[r] & [U'/R']
}
$$
は $2$-ファイバー積正方形である。
\end{lemma}

\begin{proof}
逆同型 $i:R\to R$ および $i':R'\to R'$ を補題の主張の（第1の）カルテシアン図式に適用すると、
$$
\xymatrix{
R \ar[d]_t \ar[r]_f & R' \ar[d]^{t'} \\
U \ar[r]^f & U'
}
$$
もカルテシアンであることが分かる。Lemma \ref{spaces-groupoids-lemma-cartesian-square-of-morphism} により、
$2$-ファイバー正方形
$$
\xymatrix{
[U''/R''] \ar[d] \ar[r] & [U/R] \ar[d] \\
\mathcal{S}_{U'} \ar[r] & [U'/R']
}
$$
がある。ここで $U''=U\times_{f,U',t'}R'$ および
$R''=R\times_{f\circ s,U',t'}R'$ である。上により
$(t,f):R\to U''$ は同型であり、また
$$
R'' =
R \times_{f \circ s, U', t'} R' =
R \times_{s, U} U \times_{f, U', t'} R' =
% P09-ERR-1335
R \times_{s, U, t} \times R.
$$
である。明示的には、同型 $R\times_{s,U,t}R\to R''$ は規則
$(r_0,r_1)\mapsto(r_0,f(r_1))$ によって与えられる。さらに、$s'',t'',c''$ は写像
$$
R \times_{s, U, t} R \to R,
\quad
s''(r_0, r_1) = r_1, \quad t''(r_0, r_1) = c(r_0, r_1)
$$
および
$$
\begin{matrix}
c'' : &
(R \times_{s, U, t} R) \times_{s'', R, t''} (R \times_{s, U, t} R)
&
\longrightarrow
&
R \times_{s, U, t} R, \\
&
((r_0, r_1), (r_2, r_3)) &
\longmapsto & (c(r_0, r_2), r_3).
\end{matrix}
$$
に移る。同型
$$
R \times_{s, U, s} R \longrightarrow R \times_{s, U, t} R,
\quad
(r_0, r_1) \longmapsto (c(r_0, i(r_1)), r_1)
$$
を前合成すると、$t''$ および $s''$ は $\text{pr}_0$ および $\text{pr}_1$ に移り、
$c''$ は
$\text{pr}_{02}:R\times_{s,U,s}R\times_{s,U,s}R\to R\times_{s,U,s}R$ に移る。
したがって同型 $[U''/R'']\cong[R/R\times_{s,U,s}R]$ が存在することが分かる。ここで
代数空間における亜群 $(R,R\times_{s,U,s}R,s'',t'',c'')$ は、自明な亜群
$(U,U,\text{id},\text{id},\text{id})$ の $s:R\to U$ を介した制限である。
$s:R\to U$ は右逆をもつので fppf 層の全射である。したがって射
$$
[U''/R''] \cong [R/R \times_{s, U, s} R]
\longrightarrow
[U/U] = \mathcal{S}_U
$$
は Lemma \ref{spaces-groupoids-lemma-quotient-stack-restrict-equivalence} により同値である。これで補題が示された。
\end{proof}





\section{慣性と商スタック}
\label{spaces-groupoids-section-inertia}

\noindent
亜群に値をとるスタックの（相対）慣性スタックは Stacks, Section
\ref{stacks-section-the-inertia-stack} で定義されている。ファイバー圏の設定における実際の構成と
その性質のいくつかは Categories, Section \ref{categories-section-inertia} にある。

\begin{lemma}
\label{spaces-groupoids-lemma-presentation-inertia}
$B\to S$ および $(U,R,s,t,c)$ を Definition \ref{spaces-groupoids-definition-quotient-stack} (1) のようなものとする。
$G/U$ を亜群 $(U,R,s,t,c,e,i)$ の固定化群代数空間とする。Definition
\ref{spaces-groupoids-definition-stabilizer-groupoid} を参照せよ。$R'=R\times_{s,U}G$ とおき、次のように定める。
\begin{enumerate}
\item $s':R'\to G$、$(r,g)\mapsto g$。
\item $t':R'\to G$、$(r,g)\mapsto c(r,c(g,i(r)))$。
\item $c':R'\times_{s',G,t'}R'\to R'$、
% P09-ERR-1336
$((r_1,g_1),(r_2,g_2)\mapsto(c(r_1,r_2),g_1)$。
\end{enumerate}
このとき $(G,R',s',t',c')$ は $B$ 上の代数空間における亜群であり、
$$
\mathcal{I}_{[U/R]} = [G/ R'].
$$
すなわち、付随する商スタックは $[U/R]$ の慣性スタックである。
\end{lemma}

\begin{proof}
Stacks, Lemma \ref{stacks-lemma-stackification-inertia} により、
$\mathcal{I}_{[U/_{\!p}R]}=[G/_{\!p}R']$ を証明すれば十分である。
$T$ を $S$ 上のスキームとする。慣性ファイバー圏の対象で $[U/_{\!p}R]$ のもので $T$ 上にあるものは
対 $(x,g)$ によって与えられることを思い出そう。ここで $x$ は
$[U/_{\!\!p}R]$ の対象で $T$ 上にあり、$g$ は $x$ の自己同型で、その $T$ 上のファイバー圏におけるものである。
言い換えると、$x:T\to U$ および $g:T\to R$ であって
$x=s\circ g=t\circ g$ を満たすものである。これは正確に $g:T\to G$ であることを意味する。
慣性ファイバー圏における $(x,g)\to(y,h)$ の射で $T$ 上にあるものは
$r:T\to R$ で $s(r)=x$、$t(r)=y$、および $c(r,g)=c(h,r)$ を満たすものによって与えられる。
Categories, Lemma \ref{categories-lemma-inertia-fibred-category} の可換図式を参照せよ。式で書けば
$$
h = c(r, c(g, i(r))) = c(c(r, g), i(r)).
$$
である。記法 $s(r)$ などは $s\circ r$ などの略記である。
% P09-ERR-1337
$r_1:(x_2,g_2)\to(x_1,g_1)$ と $r_2:(x_1,g_1)\to(x_2,g_2)$ の合成は
$c(r_1,r_2):(x_1,g_1)\to(x_3,g_3)$ である。

\medskip\noindent
上では、$g$ を $(x,g)$ の代わりに、$\mathcal{I}_{[U/_{\!p}R]}$ の対象で $T$ 上にあるものに対して
% P09-ERR-1338
書くこともできたことに注意せよ。実際、$x$ は $g$ の像で、構造射 $G\to U$ によるものである。
このとき射 $g\to h$ で $\mathcal{I}_{[U/_{\!p}R]}$ におけるものかつ $T$ 上にあるものは、
射 $r':T\to R'$ で $s'(r')=g$ および $t'(r')=h$ を満たすものに正確に対応する。
さらに、合成は (3) で説明した規則に対応する。これで補題が証明された。
\end{proof}

\begin{lemma}
\label{spaces-groupoids-lemma-2-cartesian-inertia}
$B\to S$ および $(U,R,s,t,c)$ を Definition \ref{spaces-groupoids-definition-quotient-stack} (1) のようなものとする。
$G/U$ を亜群 $(U,R,s,t,c,e,i)$ の固定化群代数空間とする。Definition
\ref{spaces-groupoids-definition-stabilizer-groupoid} を参照せよ。標準的 $2$-カルテシアン図式
$$
\xymatrix{
\mathcal{S}_G \ar[r] \ar[d] & \mathcal{S}_U \ar[d] \\
\mathcal{I}_{[U/R]} \ar[r] & [U/R]
}
$$
が存在し、$(\Sch/S)_{fppf}$ の亜群に値をとるスタックの図式である。
\end{lemma}

\begin{proof}
Lemma \ref{spaces-groupoids-lemma-criterion-fibre-product} により、Lemma
\ref{spaces-groupoids-lemma-presentation-inertia} の射 $s':R'\to G$ が、$s$ の基底変換で、構造射 $G\to U$ によるものと
% P09-ERR-1339
同型であることを証明すれば十分である。この基底変換性は $s'$ の構成から明らかである。
\end{proof}






\section{ジェルブと商スタック}
\label{spaces-groupoids-section-gerbes}

\noindent
この節では、商スタックを Stacks, Section \ref{stacks-section-gerbes} の議論、特に
Stacks, Definition \ref{stacks-definition-gerbe-over-stack-in-groupoids} で定義されるジェルブと関係づける。
一般に、この節に現れる亜群に値をとるスタックは代数スタックではない。

\begin{lemma}
\label{spaces-groupoids-lemma-when-gerbe}
記法と仮定を Lemma \ref{spaces-groupoids-lemma-quotient-stack-functorial} のようなものとする。商スタックの射
$$
[f] : [U/R] \longrightarrow [U'/R']
$$
は $[U/R]$ を $[U'/R']$ 上のジェルブにする。ただし $f:U\to U'$ および
$R\to R'|_U$ が fppf 層の全射である場合とする。ここで $R'|_U$ は $R'$ の
$U$ への制限であり、$f:U\to U'$ を介する。
\end{lemma}

\begin{proof}
Stacks, Lemma \ref{stacks-lemma-when-gerbe} の性質 (2)(a) および (2)(b) が成り立つことを確認する。
性質 (2)(a) は $U\to U'$ が層の全射なので成り立つ
（$[U'/R']$ の対象が局所的に $U'$ から来ることを見るには Lemma
\ref{spaces-groupoids-lemma-quotient-stack-objects} を用いよ）。(2)(b) を証明するため、$x,y$ を
$[U/R]$ の対象でスキーム $T/S$ 上にあるものとする。$x',y'$ を $x,y$ の像で圏
% P09-ERR-1340
$[U'/'R]_T$ におけるものとする。条件 (2)(b) は、層の写像
$$
\mathit{Isom}(x, y) \longrightarrow \mathit{Isom}(x', y')
$$
が $(\Sch/T)_{fppf}$ 上で全射であることの確認を要求する。これを見るには $T$ 上 fppf 局所的に
% P09-ERR-1341
作業して、$a,b\in U(T)$ から来ると仮定してよい。この場合 $x',y'$ は $f\circ a,f\circ b$ に対応する。
Lemma \ref{spaces-groupoids-lemma-quotient-stack-morphisms} により、表示した層の写像はこの場合
$$
T \times_{(a, b), U \times_B U} R
\longrightarrow
T \times_{f \circ a, f \circ b, U' \times_B U'} R' =
T \times_{(a, b), U \times_B U} R'|_U.
$$
となる。したがって $R\to R'|_U$ が $(\Sch/S)_{fppf}$ 上の fppf 層の全射であるという仮定は、
所望の全射性を含意する。
\end{proof}

\begin{lemma}
\label{spaces-groupoids-lemma-group-quotient-gerbe}
$S$ をスキームとする。$B$ を $S$ 上の代数空間とする。$G$ を $B$ 上の群代数空間とする。
$B$ に $G$ の自明な作用を入れる。射
$$
[B/G] \longrightarrow \mathcal{S}_B
$$
（Lemma \ref{spaces-groupoids-lemma-quotient-stack-arrows}）は $[B/G]$ を $B$ 上のジェルブにする。
\end{lemma}

\begin{proof}
射 $B\to B$ および $B\times_B G\to B$ は層の射として全射なので、Lemma
\ref{spaces-groupoids-lemma-when-gerbe} から直ちに従う。
\end{proof}






\section{商スタックと大サイトの変更}
\label{spaces-groupoids-section-bigger-site}

\noindent
初読の際にはこの節を飛ばすことを勧める。
スタックの引き戻しは Stacks, Section \ref{stacks-section-inverse-image}
で定義されている。

\begin{lemma}
\label{spaces-groupoids-lemma-quotient-stack-change-big-site}
大サイト $\Sch_{fppf}$ および $\Sch'_{fppf}$ が与えられているとする。
$\Sch_{fppf}$ が $\Sch'_{fppf}$ に含まれると仮定する。Topologies, Section
\ref{topologies-section-change-alpha} を参照せよ。$S \in \Ob(\Sch_{fppf})$ とする。
$B, U, R \in \Sh((\Sch/S)_{fppf})$ を代数空間とし、
$(U, R, s, t, c)$ を $B$ 上の代数空間における亜群とする。
% P09-ERR-1342
$f : (\Sch'/S)_{fppf} \to (\Sch/S)_{fppf}$ を、包含関手
$u : \Sch_{fppf} \to \Sch'_{fppf}$ に対応するサイトの射とする。
このとき、標準的同値
$$
[f^{-1}U/f^{-1}R]
\longrightarrow
f^{-1}[U/R]
$$
が存在し、これは $(\Sch'/S)_{fppf}$ 上の亜群に値をとるスタックの同値である。
\end{lemma}

\begin{proof}
Spaces, Lemma \ref{spaces-lemma-change-big-site} により
$f^{-1}B, f^{-1}U, f^{-1}R \in \Sh((\Sch'/S)_{fppf})$ は代数空間であり、したがって
$(f^{-1}U, f^{-1}R, f^{-1}s, f^{-1}t, f^{-1}c)$ は
$f^{-1}B$ 上の代数空間における亜群であることに注意せよ。ゆえに、この主張は意味をなす。

\medskip\noindent
圏 $u_p[U/_{\!p}R]$ は、圏 $u_{pp}[U/_{\!p}R]$ を射の
% P09-ERR-1343
右乗法的系 $I$ に関して局所化したものである。$u_{pp}[U/_{\!p}R]$ の対象は三つ組
$$
(T', \phi : T' \to T, x)
$$
である。ここで $T' \in \Ob((\Sch'/S)_{fppf})$、
$T \in \Ob((\Sch/S)_{fppf})$、$\phi$ は $S$ 上のスキームの射であり、
$x : T \to U$ は $(\Sch/S)_{fppf}$ 上の層の射である。
スキームの射 $\phi : T' \to T$ を与えることは、射
$\phi : T' \to u(T)$ を与えることと同じであることに注意せよ。また $u(T)$ は
$f^{-1}T$ を表現するので、これは射 $T' \to f^{-1}T$ を与えることとも同じである。
さらに Spaces, Lemma \ref{spaces-lemma-fully-faithful} により、代数空間上の
$f^{-1}$ は充満忠実なので、$x$ を射 $x : f^{-1}T \to f^{-1}U$ と考えることもできる。
以後、このような同一視は断りなく行う。射
$$
(a, a', \alpha) :
(T'_1, \phi_1 : T'_1 \to T_1, x_1)
\longrightarrow
(T'_2, \phi_2 : T'_2 \to T_2, x_2)
$$
で $u_{pp}[U/_{\!p}R]$ のものは可換図式
$$
\xymatrix{
& & U \\
T'_1 \ar[d]_{a'} \ar[r]_{\phi_1} &
T_1 \ar[d]_a \ar[ru]^{x_1} \ar[r]_\alpha &
R \ar[d]^t \ar[u]_s \\
T'_2 \ar[r]^{\phi_2} &
T_2 \ar[r]^{x_2} &
U
}
$$
である。また、このような射が $I$ の元であるための必要十分条件は
$T'_1 = T'_2$ かつ $a' = \text{id}$ であることである。関手
$$
u_{pp}[U/_{\!p}R] \longrightarrow [f^{-1}U/_{\!p}f^{-1}R]
$$
を、対象上では
$$
(T', \phi : T' \to T, x) \longmapsto (x \circ \phi : T' \to f^{-1}U)
$$
により、上のような射上では
$$
(a, a', \alpha) \longmapsto (\alpha \circ \phi_1 : T'_1 \to f^{-1}R)
$$
により定義する。$I$ の元が同型射へ写されることは明らかである。実際、
$(f^{-1}U, f^{-1}R, f^{-1}s, f^{-1}t, f^{-1}c)$ は
$f^{-1}B$ 上の代数空間における亜群である。
したがって、この関手は標準的な仕方で関手
$$
u_p[U/_{\!p}R] \longrightarrow [f^{-1}U/_{\!p}f^{-1}R]
$$
を経由する。スタック化を施すと、スタックの関手
$$
f^{-1}[U/R] \longrightarrow [f^{-1}U/f^{-1}R]
$$
を得る。これは $(\Sch'/S)_{fppf}$ 上の関手である。実際、Stacks, Lemma
\ref{stacks-lemma-technical-up} により、
スタック $f^{-1}[U/R]$ は $u_p[U/_{\!p}R]$ のスタック化である。

\medskip\noindent
ここまででスタックの射が得られた。この射が同値であることを確かめるには、これが充満忠実であり、
かつ対象が局所的に本質的像に入ることを示せば十分である。Stacks, Lemmas
\ref{stacks-lemma-characterize-ff} および
\ref{stacks-lemma-characterize-essentially-surjective-when-ff} を参照せよ。
% P09-ERR-1344
対象に関する主張は、$f^{-1}R$ が全射エタール射
$f^{-1}W \to f^{-1}R$ をもつことから従う。ここで $W$ は
$(\Sch/S)_{fppf}$ のある対象である。関手が「充満」であることを示すには、
射が局所的に関手の像に入ることを示せば十分であり、これは $f^{-1}U$ が全射エタール射
$f^{-1}W \to f^{-1}U$ をもつことから従う。ここでも $W$ は $(\Sch/S)_{fppf}$ のある対象である。
関手が忠実であることの証明は省略する。
\end{proof}










\section{分離性の条件}
\label{spaces-groupoids-section-separation}

\noindent
これは実際には、射 $j : R \to U \times_B U$ に対する条件を意味する。ただし、
代数空間における亜群 $(U, R, s, t, c)$ が $B$ 上で与えられているものとする。
前節と同じく、まず対応する図式を定式化する。

\begin{lemma}
\label{spaces-groupoids-lemma-diagram-diagonal}
$B \to S$ を Section \ref{spaces-groupoids-section-notation} のようなものとする。
$(U, R, s, t, c)$ を $B$ 上の代数空間における亜群とする。
$G \to U$ を固定化群代数空間とする。
% P09-ERR-1345
次の可換図式において
$$
\xymatrix{
R \ar[d]^{\Delta_{R/U \times_B U}} \ar[rrr]_{f \mapsto (f, s(f))} & & &
R \times_{s, U} U \ar[d] \ar[r] & U \ar[d] \\
R \times_{(U \times_B U)} R \ar[rrr]^{(f, g) \mapsto (f, f^{-1} \circ g)} & & &
R \times_{s, U} G \ar[r] & G
}
$$
左側の二つの水平射は同型射であり、右側の正方形はファイバー積の正方形である。
\end{lemma}

\begin{proof}
省略する。
定義と、代数幾何における関手的な点の見方に関する演習である。
\end{proof}

\begin{lemma}
\label{spaces-groupoids-lemma-diagonal}
$B \to S$ を Section \ref{spaces-groupoids-section-notation} のようなものとする。
$(U, R, s, t, c)$ を $B$ 上の代数空間における亜群とする。
$G \to U$ を固定化群代数空間とする。
\begin{enumerate}
\item 次の条件は同値である。
\begin{enumerate}
\item $j : R \to U \times_B U$ は分離的である。
\item $G \to U$ は分離的である。
\item $e : U \to G$ は閉埋め込みである。
\end{enumerate}
\item 次の条件は同値である。
\begin{enumerate}
\item $j : R \to U \times_B U$ は局所分離的である。
\item $G \to U$ は局所分離的である。
\item $e : U \to G$ は埋め込みである。
\end{enumerate}
\item 次の条件は同値である。
\begin{enumerate}
\item $j : R \to U \times_B U$ は準分離的である。
\item $G \to U$ は準分離的である。
\item $e : U \to G$ は準コンパクトである。
\end{enumerate}
\end{enumerate}
\end{lemma}

\begin{proof}
群代数空間 $G \to U$ は $R \to U \times_B U$ の基底変換であり、この基底変換は
対角射 $U \to U \times_B U$ による。Lemma \ref{spaces-groupoids-lemma-groupoid-stabilizer} を参照せよ。
したがって $j$ が分離的（それぞれ局所分離的、準分離的）ならば、$G \to U$ は
分離的（それぞれ局所分離的、準分離的）である。Morphisms of Spaces, Lemma
\ref{spaces-morphisms-lemma-base-change-separated} を参照せよ。
ゆえに (1)、(2)、(3) のそれぞれで (a) $\Rightarrow$ (b) が成り立つ。

\medskip\noindent
逆に $G \to U$ が分離的（それぞれ局所分離的、準分離的）ならば、射
$e : U \to G$ は構造射 $G \to U$ の切断なので、閉埋め込み（それぞれ埋め込み、準コンパクト）である。
Morphisms of Spaces, Lemma \ref{spaces-morphisms-lemma-section-immersion} を参照せよ。
ゆえに (1)、(2)、(3) のそれぞれで (b) $\Rightarrow$ (c) が成り立つ。

\medskip\noindent
$e$ が閉埋め込み（それぞれ埋め込み、準コンパクト）ならば、Lemma
\ref{spaces-groupoids-lemma-diagram-diagonal} の結果（および Spaces, Lemma
\ref{spaces-lemma-base-change-immersions} と Morphisms of Spaces, Lemma
\ref{spaces-morphisms-lemma-base-change-quasi-compact}）により、
$\Delta_{R/U \times_B U}$ は閉埋め込み（それぞれ埋め込み、準コンパクト）である。
ゆえに (1)、(2)、(3) のそれぞれで (c) $\Rightarrow$ (a) が成り立つ。
\end{proof}
